% ============================================================
% 第10章：人工智能在医工融合中的应用
% ============================================================

\setcounter{chapter}{9}
\chapter{人工智能在医工融合中的应用}

在医生工作站前，人工智能可以辅助核对医嘱中的药物、检查和执行时间；患者离开医院后，手腕上的智能设备仍在记录心率、血氧和活动变化；康复训练室里，传感器捕捉患者每一次抬手、迈步和重心转移；中医方剂分析平台中，药物之间的配伍关系被整理成可查询的知识网络。人工智能不再只是影像科里的识别模型，而是逐渐进入诊疗决策、院外监测、康复训练和医学知识整理等具体环节，成为医工融合系统中的重要组成部分。

第9章介绍了医学数据的来源、类型和基本应用。本章将在此基础上，把视角进一步转向具体场景，重点观察人工智能如何与医学知识、工程系统和临床流程结合起来。本章围绕四个典型应用展开，包括医嘱智能判别与规范评分、智能穿戴监测与健康评估、智能康复系统与功能训练，以及中药方剂解读与组方分析。通过这些内容，读者能够进一步理解人工智能如何从数据分析走向实际应用，并在医院、家庭、康复室和中医诊疗等场景中发挥作用。

\section{医嘱智能判别与规范评分}

夜班病房里，医生刚提交一组新医嘱。屏幕上只是几行文字，实际却牵连着药物选择、检查安排、护理要求、监测项目和执行时间。某项药物是否适合当前诊断，必要检查是否已经开立，监测项目是否跟上病情变化，治疗顺序是否符合临床要求，这些内容都需要逐项核对。医嘱数量少时，医生可以依靠经验仔细检查；当医嘱数量较多、内容较长、患者情况较复杂时，单纯依靠人工复核就容易耗费大量时间。

医嘱智能判别与规范评分，正是针对这一类临床需求建立的人工智能应用。系统首先接收医生开立的医嘱文本，并从中识别诊断、检查、药物、治疗、护理、监测和时间等关键信息；随后调用规则库、金标准医嘱和知识图谱，对医嘱内容进行逐项比对；最后生成差异提示、规范评分和解释性结果。这样，原本连续书写的医嘱文本可以被整理为结构清楚、依据明确、便于复核的审核结果。

如图10.1所示，医嘱智能判别与规范评分通常包括医嘱数据整理、规则提取、知识图谱构建、智能审核与评分反馈等环节。前端负责接收和解析医嘱，后端负责调用规则、图谱和模型完成判别，最终向医生输出结构化、可解释的审核结果。通过这一过程，人工智能能够参与医嘱质量管理中的信息整理、标准对照和结果提示，为临床医嘱复核提供辅助支持。

\begin{figure}[H]
\centering
\begin{tcolorbox}[width=0.88\textwidth, colback=gray!10, colframe=gray!50, boxrule=0.5pt, arc=2pt]
\centering
\vspace{30pt}
{\large\color{gray!70} [图片占位]}
\vspace{10pt}
{\small\color{gray!50} 图10.1 医嘱智能判别与规范评分总体框架示意图}
\vspace{25pt}
\end{tcolorbox}
\caption*{\footnotesize\itshape\color{gray!50} 图10.1 医嘱智能判别与规范评分总体框架示意图}
\end{figure}

\subsection{医嘱数据整理与规则提取}

一条医嘱在医生眼中往往很清楚，在计算机眼中却可能是一段拥挤的文字。药品名称、剂量、频次、给药途径、检查项目、护理要求和执行时间可能混在同一句话里。要让系统审核医嘱，第一步不是直接调用模型，而是先把医嘱拆开，让机器知道每个词语对应哪类医学含义。

如图10.2所示，医嘱数据整理与规则提取通常包括医嘱要素拆解、结构化抽取、金标准对比与规则归纳等环节。整理的目的不是把临床工作变复杂，而是把医生自然书写的医嘱转化为机器可以读取、比较和追溯的数据结构。

\begin{figure}[H]
\centering
\begin{tcolorbox}[width=0.88\textwidth, colback=gray!10, colframe=gray!50, boxrule=0.5pt, arc=2pt]
\centering
\vspace{30pt}
{\large\color{gray!70} [图片占位]}
\vspace{10pt}
{\small\color{gray!50} 图10.2 医嘱数据整理与规则提取流程示意图}
\vspace{25pt}
\end{tcolorbox}
\caption*{\footnotesize\itshape\color{gray!50} 图10.2 医嘱数据整理与规则提取流程示意图}
\end{figure}

\subsubsection*{\textbf{1. 医嘱的基本类型与构成要素}}

住院医嘱就像患者住院期间的“诊疗任务单”。医生写下医嘱后，药房、护士站、检验科、影像科和治疗团队都会根据它开展后续工作。按照执行特点，医嘱通常分为长期医嘱和临时医嘱。长期医嘱像住院治疗的主线，负责较长时间内的用药、护理、饮食、监测和复查安排；临时医嘱则更像应对突发情况的即时指令，用于临时检查、紧急用药、会诊申请或单次处置。

一条医嘱虽然看起来只有几行字，但里面至少要交代清楚四件事：给谁做、做什么、怎么做、什么时候做。对应到医嘱数据中，就是诊疗对象信息、诊疗项目信息、执行参数信息和时间状态信息。患者身份、诊断、病区和床号说明医嘱面向谁；药品、检查、治疗和护理说明要做什么；剂量、频次、途径和部位说明怎么做；开立时间、执行时间和执行状态说明什么时候做、是否完成。

医嘱要素是否完整，直接影响后续执行和审核。若药物医嘱缺少剂量或频次，执行人员就难以准确给药；若检查医嘱缺少部位或时间要求，检查安排和结果判读可能受到影响；若治疗或护理医嘱缺少执行时间，系统也难以判断其是否符合诊疗阶段。因此，医嘱数据整理首先需要明确对象、项目、参数和时间状态等基本要素，为后续结构化抽取、规则比对和规范评分提供可靠基础。

医嘱数据的结构化抽取

在医院信息系统中，医嘱的写法并不总是完全一致。有些医嘱来自标准模板，药品、剂量、频次和途径等字段已经相对清楚；有些医嘱则由医生根据患者情况直接录入，可能把检查要求、治疗安排、监测时间和补充说明写在同一段文字中。结构化抽取的任务，就是从这些不同形式的医嘱中识别关键信息，并将其整理为统一字段，使系统能够进一步完成比对、审核和评分。

结构化抽取通常包括分词、术语识别、要素归类和关系抽取四个步骤。分词是将连续文本切分为具有医学含义的片段；术语识别用于找出诊断、药品、检查、治疗、护理、监测指标和时间表达；要素归类是将识别出的内容放入对应字段；关系抽取则进一步确定不同要素之间的搭配关系，如药品对应的剂量和频次、检查对应的部位、监测对应的时间点等。

在实际应用中，结构化抽取通常需要多种方法配合完成。规则词典适合识别固定术语，具有稳定、可解释的特点；传统机器学习方法可以处理一定程度的表达变化；预训练语言模型更适合处理长句、缩写和自由文本中的复杂语义。多种方法结合使用，可以提高医嘱要素识别的准确性和适应性，为后续规则比对、知识图谱构建和规范评分提供基础。

\subsubsection*{\textbf{3. 金标准医嘱与对比依据}}

医生复核医嘱时，通常不会只看某条医嘱写得像不像，而是会把它放回具体疾病和治疗阶段中判断。患者处于入院评估期、急性治疗期还是康复随访期，需要完成的检查、治疗、护理和监测安排并不相同。系统进行医嘱审核时，也需要这样一套明确的参照标准。金标准医嘱就是这一参照，它是指在某一疾病、某一治疗阶段或某一特定情形下，依据指南、临床路径、质控标准或专家共识整理形成的规范化医嘱集合。

金标准医嘱通常包括适用诊断、治疗阶段、医嘱类别、医嘱项目、具体药物或措施、控制目标、触发条件、禁忌事项和依据来源等内容。通过这些信息，系统不仅可以判断某项医嘱是否出现，还可以进一步判断它适用于哪个阶段、是否满足触发条件、是否需要配套监测，以及是否存在禁忌或不推荐使用的情况。

明确的对比依据能够使审核结果更容易解释。当系统提示某项检查需要核对时，应当说明该检查对应的疾病诊断、诊疗阶段或评估目的；当系统提示某类药物需要监测时，也应能够追溯到药物与监测指标之间的关系。这样生成的审核结果不再只是简单提示“可能不规范”，而是能够给出依据明确、指向清楚的辅助说明。

\subsubsection*{\textbf{4. 规则提取与表达方式}}

临床医生阅读指南时，能够根据经验理解其中的含义。指南里写着“必要时复查影像”、“用药期间注意监测电解质”，医生知道这意味着什么情境下该做什么。但对计算机来说，这些表述还不够明确。系统不能靠经验领会文字背后的临床意图，它需要清楚的条件、对象和结果。规则提取的任务，就是把这些医学要求转化为机器能够执行的判断规则。

规则提取通常包括语义识别、条件归纳、冲突消解和层级组织等步骤。语义识别用于找出诊断、药物、检查、治疗和监测等关键信息；条件归纳用于把临床描述整理成“在什么情况下需要做什么”的判断条件；冲突消解用于处理不同指南、路径或专家共识之间可能存在的差异；层级组织则按照疾病类型、诊疗阶段、医嘱类别和适用范围进行管理，便于系统后续调用和维护。

规则的表达方式可以根据任务复杂程度进行选择。对于项目明确、逻辑简单的内容，可以采用关键词匹配和模板规则，用于判断某项检查、药物或监测是否出现；对于涉及同义表达、上下位概念、条件触发和阶段顺序的内容，则更适合采用本体、语义网或知识图谱进行表达。医嘱审核通常需要多种规则形式配合，才能覆盖简单漏项、用药监测、禁忌判断和复杂诊疗流程等不同问题。

\subsubsection*{\textbf{5. 小结}}

本节介绍了医嘱数据整理与规则提取的基本过程。系统需要先识别医嘱中的诊疗对象、医嘱项目、执行参数和时间状态等要素，再通过结构化抽取将自由文本转化为可计算的数据形式。在此基础上，金标准医嘱为系统提供对比依据，规则提取则将指南、临床路径和专家经验整理为可调用的审核条件。

\subsection{知识图谱构建与存储管理}

在医嘱审核中，系统面对的不是一张简单清单，而是一组相互关联的医学知识。疾病诊断会影响检查项目的选择，治疗阶段会决定用药和护理安排，药物使用又会带出相应的监测要求和禁忌判断。例如，高血压患者是否需要监测肾功能，脑出血患者何时复查影像，某类药物使用后是否需要关注电解质变化，这些问题都不能只看单个医嘱项目本身，而要放在疾病、阶段、治疗和监测之间的关系中理解。

知识图谱能够把这些分散的医学知识组织成由节点和关系构成的网络。节点可以表示疾病、药物、检查项目、治疗措施、护理要求、监测指标和诊疗阶段；关系则用来表示它们之间的联系。它将指南、临床路径、金标准医嘱和专家经验中的内容进行整理，形成可查询、可推理、可解释的知识结构。当系统接收到一条实际医嘱时，可以根据患者诊断和医嘱内容，在图谱中查找相关标准要求，判断是否存在检查遗漏、监测不足、用药禁忌或阶段不匹配等情况，并为审核结果提供依据。

如图10.3所示，医嘱知识图谱构建与存储管理通常包括知识来源整理、实体与关系抽取、知识融合、图数据库存储和持续维护等环节。知识来源整理负责汇集指南、共识、药品说明书和金标准医嘱；实体与关系抽取负责识别疾病、药物、检查和监测等核心内容及其联系；知识融合用于处理不同来源之间的名称差异和规则冲突；图数据库存储则负责保存节点、关系和属性；持续维护用于根据指南更新、临床需求变化和系统反馈不断修订图谱内容。

\begin{figure}[H]
\centering
\begin{tcolorbox}[width=0.88\textwidth, colback=gray!10, colframe=gray!50, boxrule=0.5pt, arc=2pt]
\centering
\vspace{30pt}
{\large\color{gray!70} [图片占位]}
\vspace{10pt}
{\small\color{gray!50} 图10.3 医嘱知识图谱构建与存储管理流程示意图}
\vspace{25pt}
\end{tcolorbox}
\caption*{\footnotesize\itshape\color{gray!50} 图10.3 医嘱知识图谱构建与存储管理流程示意图}
\end{figure}

\subsubsection*{\textbf{1. 实体识别与关系抽取}}

一张医嘱知识图谱的形成，有点像把散落在病历、指南和医嘱中的医学信息重新画成一张关系图。系统首先要找到图里的“点”，也就是医学实体。疾病名称、药品名称、检查项目、治疗措施、护理项目、监测指标和时间阶段等，都可以作为实体进入图谱。例如，在脑血管疾病相关医嘱中，头颅CT、GCS评分、甘露醇、血常规、电解质和颅内压监测等内容，都可以被识别为图谱中的节点。

实体识别之后，还需要进一步确定这些节点之间的关系。诊断可以与不同诊疗阶段建立对应关系，阶段可以与应执行的检查、治疗和护理项目建立关联，药物可以与需要监测的指标建立联系，特殊病情或适应证也可以触发相应的医嘱项目。通过关系抽取，系统能够知道某个节点不是孤立存在的，而是处在疾病、阶段、治疗和监测共同构成的医学网络中。

医学实体识别的难点在于表达方式并不统一。同一项检查可能被写作CT、头颅CT或头颅CT平扫；肝肾功能可能对应多个实验室指标；抗感染治疗则可能涉及具体药物、适应证和病原学判断。图谱构建需要对这些同义词、缩写和不同粒度的表达进行规范化处理，将它们映射到明确的节点上。这样，系统在后续审核时才能减少名称差异带来的干扰，使医嘱比对和规则推理更加稳定。

\subsubsection*{\textbf{2. 知识融合与图谱构建}}

在医院里，同一个诊疗要求可能出现在不同地方。诊疗指南会给出原则性建议，临床路径会安排具体步骤，药品说明书会强调用法用量和禁忌，质控标准会关注哪些项目必须核对，医生经验则常常补充实际执行中的细节。这些资料都很重要，但它们并不会自动排成整齐的一套知识。知识融合要做的，就是把来自不同来源的信息放到同一套概念和关系框架中，使系统能够按照统一标准进行查询、比对和推理。

知识融合通常包括实体对齐、关系归并和冲突消解。实体对齐用于把不同名称映射到同一概念，例如将不同写法的疾病名称、检查项目或药品名称统一到同一节点。关系归并用于整理含义相近的关系表达，例如将建议检查、应完成检查、需复查等内容归入相应的检查关系。冲突消解则用于处理不同来源之间可能存在的差异，如推荐时间、用药范围或监测要求不完全一致的情况。

图谱构建还需要考虑临床流程的实际特点。医嘱之间并不总是按照一条直线依次发生。有些内容具有明确先后顺序，如入院评估、急性期处理和康复管理；有些内容可以并行开展，如监测、护理和部分实验室检查；还有一些内容只在特定病情变化或特殊条件下触发。为了表达这种关系，系统可以采用多通道偏序结构。主流程保留必要的阶段顺序，分支内容通过触发条件、依赖关系和适用范围进行表达，并行内容则不强行排列先后。这样既能保留关键时间逻辑，也能避免把并行任务误判为顺序错误。

经过知识融合与图谱构建后，原本分散在不同资料中的医学知识被组织成较为稳定的知识网络。系统不仅能够查询某一疾病对应的检查、治疗和监测要求，也能够进一步判断这些项目属于哪个阶段、依赖什么条件、是否与其他医嘱存在关联，为后续医嘱审核提供了更清晰、可追溯的医学依据。

\subsubsection*{\textbf{3. 图数据库存储与查询}}

知识图谱建成以后，需要被放入能够支持关系查询的数据库中。普通表格像一本按项目排列的登记册，能够记录每个医嘱项目的名称和说明；图数据库则像一张展开的诊疗关系图，疾病、阶段、药物、检查、治疗、护理和监测之间都有清晰的连接。图数据库以节点和关系为基本存储单位。诊断、阶段、医嘱项目、药物、检查、监测指标和时间窗等内容可以作为节点保存，节点之间通过关系连接起来。系统查询时，可以从诊断节点找到对应治疗阶段，从阶段节点找到应执行的检查和治疗项目，从药物节点找到需要配套监测的指标，也可以从某项检查继续追溯其所属阶段和标准时间窗。

在医嘱审核中，图数据库主要支持三类查询。第一类是关系查询，用于查找疾病、药物、检查和监测之间的直接联系。第二类是路径查询，用于判断某项医嘱是否能够通过阶段、条件或适用范围与当前诊断建立关联。第三类是子图匹配，用于将实际医嘱形成的局部结构与金标准医嘱结构进行比较，从而发现覆盖情况和差异内容。

图数据库的优势在于查询过程清楚，结果也更容易解释。当系统提示某项检查需要核对时，可以追溯到对应的诊断节点、阶段节点和标准项目节点；当系统判断某项医嘱时间需要关注时，也可以找到该项目所属阶段及其时间窗。这样生成的审核结果不只是一个孤立提示，而是带有来源、路径和依据的结构化说明，更便于医生复核和后续质量管理。

\subsubsection*{\textbf{4. 知识图谱的维护与更新}}

医院里的诊疗知识并不是一本写完就不再翻动的旧书。新的指南会发布，新的药物会进入临床，检查项目和质控要求也会随着实际工作不断调整。知识图谱如果长期不更新，就像一张没有改过路线的医院导航图，表面上仍然完整，但其中某些路径可能已经不适合当前诊疗要求。

知识图谱的维护通常包括内容更新、版本管理和质量审核。内容更新用于补充新增疾病、药物、检查项目、治疗措施和审核规则；版本管理用于记录每次更新的来源、时间、修改内容和影响范围；质量审核则用于检查节点是否重复、关系是否错误、时间窗是否合理，以及不同规则之间是否存在冲突。对于临床应用而言，版本记录尤其重要，因为同一条医嘱规则在不同时间可能对应不同的指南依据。

图谱更新一般包括周期性更新和触发性更新两种方式。周期性更新适合对常见病种、常用医嘱和基础规则进行定期整理，保证图谱整体内容保持稳定；触发性更新则用于应对新指南发布、新质控要求提出、新药物使用或临床专家补充标准等情况。两种方式结合，可以使知识图谱既保持结构稳定，又能够及时吸收新的医学知识，为医嘱审核系统提供持续可靠的依据。

\subsubsection*{\textbf{5. 小结}}

本节介绍了医嘱知识图谱构建与存储管理的基本过程。系统首先通过实体识别提取疾病、药物、检查、治疗、护理、监测和时间阶段等医学对象，再通过关系抽取建立它们之间的诊疗阶段、执行要求、监测需求和触发条件等联系。在此基础上，知识融合用于统一不同来源中的名称、关系和规则表达，减少同义词、缩写和来源差异带来的影响。

知识图谱建成后，需要通过图数据库进行存储和查询，使系统能够根据诊断、阶段、医嘱项目和时间窗快速追溯标准依据。同时，知识图谱还需要持续维护和更新，以适应指南修订、新药应用、质控要求变化和临床标准补充。由此可见，知识图谱不仅是医嘱审核系统的知识存储方式，也是连接医学规范、工程系统和模型判别的重要基础。

\subsection{案例：基于大模型的医嘱审核与规范评分}

一次系统测试中，研究人员把一段脑出血相关医嘱输入审核界面。点击“开始审核”后，系统并不会简单查找关键词，而是先把文本中的诊断、药物、检查、治疗、护理、监测和时间信息逐项识别出来，再与本地图数据库中的金标准医嘱进行对照。几秒钟后，页面生成了完整的审核结果：哪些内容已经识别，哪些检查可能遗漏，哪些时间安排需要注意，最终规范评分是多少。这个过程展示了一套可以运行、可以输出、可以被医生复核的医嘱判别流程。

本案例采用二附院提供的金标准医嘱作为对照依据。金标准医嘱覆盖脑出血相关疾病在不同阶段应关注的检查、治疗、监测、护理和随访要求。系统将这些金标准内容保存到本地图数据库中，并调用本地部署的7B模型对输入医嘱进行判别与规范评分。审核完成后，系统输出实际医嘱识别结果、金标准对比结果、差异说明、时间一致性检查和综合评分，为医师复核提供依据。

\begin{figure}[H]
\centering
\begin{tcolorbox}[width=0.88\textwidth, colback=gray!10, colframe=gray!50, boxrule=0.5pt, arc=2pt]
\centering
\vspace{30pt}
{\large\color{gray!70} [图片占位]}
\vspace{10pt}
{\small\color{gray!50} 图10.4 基于大模型的医嘱智能审核系统总体架构示意图}
\vspace{25pt}
\end{tcolorbox}
\caption*{\footnotesize\itshape\color{gray!50} 图10.4 基于大模型的医嘱智能审核系统总体架构示意图}
\end{figure}

\subsubsection*{\textbf{1. 案例数据与系统设置}}

本案例采用二附院提供的标准医嘱数据作为金标准，用于判断输入医嘱是否有效、是否存在错误，以及是否可能导致误诊或漏诊。本次案例新增了24个金标准数据，覆盖高钾血症、高血压性脑出血和动脉瘤性蛛网膜下腔出血等病种。每条金标准医嘱不仅记录具体医嘱项目，还包含所属阶段、适用条件、时间要求和相关依赖关系，为后续医嘱比对和规范评分提供依据。

从内容构成看，金标准医嘱主要包括诊断、检查、治疗、护理、监测和随访等信息。诊断信息用于确定当前医嘱适用的疾病背景；检查信息用于判断必要检验和影像评估是否完整；治疗信息用于核对药物、操作和治疗方案是否合理；护理和监测信息用于判断患者状态观察是否充分；随访信息则用于评价出院后复查和长期管理安排是否完善。通过上述内容，系统能够从多个维度评价实际医嘱与标准医嘱之间的差异。

在系统设置上，本案例采用本地部署的7B模型与本地图数据库协同完成医嘱审核。7B模型主要负责自然语言理解、医嘱内容识别、差异描述和结果组织；图数据库主要负责存储金标准医嘱、阶段关系、时间窗和依赖关系，并为模型判别提供可追溯的医学依据。

\subsubsection*{\textbf{2. 金标准知识图谱构建}}

金标准医嘱中的内容并不是一组孤立条目。诊断、阶段、检查、治疗、护理、监测和随访之间往往存在明确的医学关联。若只用普通表格保存，每一条医嘱虽然能够被记录下来，但阶段归属、时间要求、触发条件和依赖关系不容易被系统直接识别。因此，本案例将金标准医嘱构建为知识图谱，使标准内容能够以节点和关系的形式被存储、查询和推理。

在知识图谱中，系统将金标准医嘱拆分为多个核心实体，并分别建立对应节点。诊断节点用于表示疾病背景，阶段节点用于表示诊疗过程中的不同时间或病程阶段，项目节点用于表示具体医嘱内容，包括药物、检查、治疗操作、护理措施、监测项目和随访安排。这样处理后，原本写在表格中的标准医嘱被转化为一组可以被系统识别和调用的知识单元，为后续医嘱比对、时间判断和规范评分提供基础。

\begin{figure}[H]
\centering
\begin{tcolorbox}[width=0.88\textwidth, colback=gray!10, colframe=gray!50, boxrule=0.5pt, arc=2pt]
\centering
\vspace{30pt}
{\large\color{gray!70} [图片占位]}
\vspace{10pt}
{\small\color{gray!50} 图10.4-1展示了金标准医嘱知识图谱中的阶段节点与关系示意图。图中以“创伤性硬膜下出血”为中心，可以看到其与诊断确认、T0急性期、T0+6-24小时、T0+1-7天、康复期和长期管理等阶段节点相连。不同阶段再进一步连接对应的计划项目，例如康复期关联慢性硬膜下血肿处置，T0+1-7天阶段关联术后管理、保守治疗等内容。这说明金标准在图数据库中并不是按照单一时间线机械排列，而是根据病种、阶段、计划项目和适用范围组织成具有层次关系的知识结构。}
\vspace{25pt}
\end{tcolorbox}
\caption*{\footnotesize\itshape\color{gray!50} 图10.4-1展示了金标准医嘱知识图谱中的阶段节点与关系示意图。图中以“创伤性硬膜下出血”为中心，可以看到其与诊断确认、T0急性期、T0+6-24小时、T0+1-7天、康复期和长期管理等阶段节点相连。不同阶段再进一步连接对应的计划项目，例如康复期关联慢性硬膜下血肿处置，T0+1-7天阶段关联术后管理、保守治疗等内容。这说明金标准在图数据库中并不是按照单一时间线机械排列，而是根据病种、阶段、计划项目和适用范围组织成具有层次关系的知识结构。}
\end{figure}

右侧节点详情显示，长期管理节点被标记为Phase（阶段）类型，并包含name（名称）、phase_type（阶段类型）、sequence（顺序号）和is_mainline（是否属于主流程）等属性。其中，is_mainline为true，说明该节点属于主流程阶段；sequence用于表示其在主流程中的相对顺序。通过这些属性，系统能够区分主流程阶段、分支阶段和并行项目，为后续判断某项医嘱是否属于对应阶段、是否符合时间要求、是否应与其他医嘱同时出现提供依据。

\begin{figure}[H]
\centering
\begin{tcolorbox}[width=0.88\textwidth, colback=gray!10, colframe=gray!50, boxrule=0.5pt, arc=2pt]
\centering
\vspace{30pt}
{\large\color{gray!70} [图片占位]}
\vspace{10pt}
{\small\color{gray!50} 图10.4-1 金标准医嘱知识图谱中的阶段节点与关系示意图}
\vspace{25pt}
\end{tcolorbox}
\caption*{\footnotesize\itshape\color{gray!50} 图10.4-1 金标准医嘱知识图谱中的阶段节点与关系示意图}
\end{figure}

\subsubsection*{\textbf{3. 医嘱下达时间判定}}

本次系统更新增加了对医嘱下达时间的判定。医嘱审核不仅要判断某项检查、治疗或监测是否出现，还要判断其是否出现在合适的诊疗阶段。对于脑出血相关医嘱而言，头颅CT复查、止血治疗、降颅压处理、感染控制、护理监测和出院随访等内容都有相应的时间要求。若只判断“有没有”，容易忽略医嘱下达过早或过晚带来的临床风险。

在时间判定中，系统首先识别输入医嘱中的时间描述，如T0、T0+2-4h、2小时后、24小时内等，并将其统一换算为小时区间。随后，系统根据医嘱实体类型，如药物、检查、治疗、护理或监测，到知识图谱中查找对应阶段及其标准时间窗。最后，将输入时间与标准时间窗进行比较：早于窗口判为提前，晚于窗口判为超时，落在窗口内判为合规，并将结果写入时间一致性检查。

\subsubsection*{\textbf{4. 医嘱识别与金标准对照}}

本案例输入的是一段脑出血相关医嘱文本，包含营养维持、预防癫痫、降低颅内压、抗感染、持续心电监测、中心静脉压监测、颅内压监测、血糖监测、静脉血栓预防、复查血气、复查脑CT以及出院后随访等多项内容。系统审核时，首先需要从这段连续文本中识别出诊断、药物、检查、治疗、护理、监测和随访等要素，再将这些要素与图数据库中召回的金标准医嘱进行对照。

\begin{figure}[H]
\centering
\begin{tcolorbox}[width=0.88\textwidth, colback=gray!10, colframe=gray!50, boxrule=0.5pt, arc=2pt]
\centering
\vspace{30pt}
{\large\color{gray!70} [图片占位]}
\vspace{10pt}
{\small\color{gray!50} 图10.4-2展示了系统界面中的输入医嘱、金标准内容和实际医嘱识别结果。界面上方为待审核医嘱文本，系统点击“开始审核”后，会根据输入内容召回对应病种的金标准医嘱。图中金标准部分显示，本案例对应的是脑室出血相关金标准，内容覆盖T0急性期、T0+24小时内、T0+1-7天、T0+1-14天和康复期等不同阶段，涉及监测、检查、评估、治疗、禁忌或不推荐药物以及长期管理要求。}
\vspace{25pt}
\end{tcolorbox}
\caption*{\footnotesize\itshape\color{gray!50} 图10.4-2展示了系统界面中的输入医嘱、金标准内容和实际医嘱识别结果。界面上方为待审核医嘱文本，系统点击“开始审核”后，会根据输入内容召回对应病种的金标准医嘱。图中金标准部分显示，本案例对应的是脑室出血相关金标准，内容覆盖T0急性期、T0+24小时内、T0+1-7天、T0+1-14天和康复期等不同阶段，涉及监测、检查、评估、治疗、禁忌或不推荐药物以及长期管理要求。}
\end{figure}

在实际医嘱识别结果中，系统将输入文本中的内容初步拆分为多个类别。其中，诊断被识别为“脑出血合并脑室出血”；药物部分识别出营养维持、预防癫痫、抗感染、降颅压、脱水治疗、止血治疗和相关用药安排；监测部分识别出持续心电监测、中心静脉压、颅内压、血糖监测以及出入量记录等内容；检查部分识别出复查血气、血常规、CRP、电解质和脑CT复查等内容。

从图10.4-2可以看出，系统不仅输出了已识别的实际医嘱，还进一步将其与金标准进行差异对照。对于已经覆盖的内容，系统能够将其归入相应类别；对于可能缺失的内容，系统会在后续差异分析中提示。例如，金标准中强调头颅非增强CT、病因鉴别、严重程度评估、脑积水评估和长期管理等内容，而实际输入医嘱虽然包含复查脑CT和多项治疗监测安排，但仍需要进一步判断是否满足金标准中对检查项目、评估项目和阶段时间的要求。

\begin{figure}[H]
\centering
\begin{tcolorbox}[width=0.88\textwidth, colback=gray!10, colframe=gray!50, boxrule=0.5pt, arc=2pt]
\centering
\vspace{30pt}
{\large\color{gray!70} [图片占位]}
\vspace{10pt}
{\small\color{gray!50} 图10.4-2 医嘱审核系统中的输入医嘱、金标准与实际医嘱识别结果}
\vspace{25pt}
\end{tcolorbox}
\caption*{\footnotesize\itshape\color{gray!50} 图10.4-2 医嘱审核系统中的输入医嘱、金标准与实际医嘱识别结果}
\end{figure}

\subsubsection*{\textbf{5. 审核结果与评分输出}}

\begin{figure}[H]
\centering
\begin{tcolorbox}[width=0.88\textwidth, colback=gray!10, colframe=gray!50, boxrule=0.5pt, arc=2pt]
\centering
\vspace{30pt}
{\large\color{gray!70} [图片占位]}
\vspace{10pt}
{\small\color{gray!50} 图10.4-3展示了本案例的差异对比与规范评分结果。系统在识别实际医嘱后，将其与召回的金标准医嘱进行逐项对照，并把审核结果分为实际医嘱识别、差异对比和评分三部分输出。绿色区域显示了系统从输入文本中识别出的诊断、药物、检查、治疗、护理和监测等内容，便于医师快速查看系统是否正确理解了原始医嘱。}
\vspace{25pt}
\end{tcolorbox}
\caption*{\footnotesize\itshape\color{gray!50} 图10.4-3展示了本案例的差异对比与规范评分结果。系统在识别实际医嘱后，将其与召回的金标准医嘱进行逐项对照，并把审核结果分为实际医嘱识别、差异对比和评分三部分输出。绿色区域显示了系统从输入文本中识别出的诊断、药物、检查、治疗、护理和监测等内容，便于医师快速查看系统是否正确理解了原始医嘱。}
\end{figure}

在差异对比部分，系统提示本案例可能存在金标准检查遗漏，主要问题为未明确包含头颅非增强CT。系统给出的说明是，CT是诊断脑室出血的重要金标准，可用于快速评估出血量、出血部位、脑积水及占位效应，并为后续治疗决策提供依据。也就是说，系统并不是简单指出“少了一项检查”，而是进一步说明该检查为什么重要，以及它在当前疾病场景中的医学意义。

在评分结果中，系统给出的总分为85/100，扣分项显示为“遗漏金标准检查：头颅非增强CT”，对应扣15分。该结果表明，审核模型能够从较长的连续医嘱文本中识别出诊断、药物、监测、检查和治疗等关键信息，并将其与金标准医嘱进行结构化对照。输入医嘱中包含营养维持、抗感染、降颅压、止血处理、心电监测、颅内压监测、血糖监测、血气复查和脑CT复查等多项内容，模型能够对这些内容进行分类整理，并进一步提示头颅非增强CT这一金标准检查项目在脑室出血诊断和病情评估中的重要作用。由此可见，该模型不仅能够完成医嘱文本识别，还能够结合金标准给出差异提示、扣分依据和医学解释，说明其在医嘱规范性审核中具有较好的实用性和可解释性。

\begin{figure}[H]
\centering
\begin{tcolorbox}[width=0.88\textwidth, colback=gray!10, colframe=gray!50, boxrule=0.5pt, arc=2pt]
\centering
\vspace{30pt}
{\large\color{gray!70} [图片占位]}
\vspace{10pt}
{\small\color{gray!50} 图10.4-3 差异对比与规范评分结果示意图}
\vspace{25pt}
\end{tcolorbox}
\caption*{\footnotesize\itshape\color{gray!50} 图10.4-3 差异对比与规范评分结果示意图}
\end{figure}

\subsubsection*{\textbf{6. 小结}}

本案例以二附院提供的金标准医嘱为对照依据，展示了基于本地图数据库和本地7B模型的医嘱审核流程。系统首先对输入医嘱进行结构化识别，提取其中的诊断、药物、检查、治疗、护理、监测和时间信息；随后调用图数据库中的金标准知识图谱，对实际医嘱与标准医嘱进行对照；最后生成差异提示、规范评分和解释性结果。本案例中，系统输出85/100的规范评分，并提示头颅非增强CT这一检查项目需要重点核对，说明模型能够围绕复杂长文本医嘱完成识别、比对、评分和结果解释。

上述审核过程体现了人工智能在医学场景中的辅助价值。模型能够将长文本医嘱中的临床信息转化为结构化内容，并结合标准知识完成快速核对和解释。对于医嘱审核这类信息密集、规范性要求较高的工作，人工智能可以承担信息整理、标准匹配和结果提示等任务，帮助临床人员更高效地把握医嘱中的关键内容，也为医学服务的规范化和智能化提供技术支持。

\section{智能穿戴监测与健康评估}

患者离开医院后，身体并不会停止变化。清晨刚起床时，血压可能悄悄升高；夜里熟睡时，血氧可能短暂下降；饭后一两个小时内，血糖曲线可能明显上扬；一次爬楼之后，心率也可能比平时恢复得更慢。这些变化大多发生在家里、路上、办公室或睡眠中，未必刚好出现在门诊测量的几分钟里。门诊检查能够记录患者到院时的状态，却很难完整呈现其日常生活中的连续变化。

智能穿戴设备的出现，使健康监测从医院检查室延伸到生活场景。佩戴在手腕上的智能手表、贴附在胸前的心电贴片、用于糖尿病管理的连续血糖监测仪，以及用于老年人安全防护的跌倒传感器，都可以持续采集人体生理和行为信息。它们记录的不再只是某一次测量结果，而是一段时间内身体状态的变化轨迹。

这些连续数据需要经过进一步整理和分析，才能真正服务于健康评估。人工智能可以对心率曲线、血氧变化、心电波形、血糖波动、睡眠片段和活动记录进行识别与判断，从中发现异常变化、评估健康状态，并将结果反馈给用户或医护人员。这样，健康管理不再完全依赖单次检查和患者回忆，而能够基于更连续、更动态的数据进行判断。

如图10.5所示，智能穿戴监测与健康评估通常包括生理参数采集、连续数据传输、健康状态识别和异常预警反馈等环节。设备端负责采集人体信号，平台端负责完成数据汇总和状态分析，医护人员或用户根据系统提示进行进一步处理。各环节相互衔接，形成从数据采集、健康评估到反馈干预的完整过程。

\begin{figure}[H]
\centering
\begin{tcolorbox}[width=0.88\textwidth, colback=gray!10, colframe=gray!50, boxrule=0.5pt, arc=2pt]
\centering
\vspace{30pt}
{\large\color{gray!70} [图片占位]}
\vspace{10pt}
{\small\color{gray!50} 图10.5 智能穿戴监测与健康评估总体框架示意图}
\vspace{25pt}
\end{tcolorbox}
\caption*{\footnotesize\itshape\color{gray!50} 图10.5 智能穿戴监测与健康评估总体框架示意图}
\end{figure}

\subsection{生理参数采集与连续监测}

生理参数采集与连续监测，是智能穿戴健康评估的基础环节。系统首先需要通过传感器获取人体生理和行为信息，常见内容包括心率、血氧、心电、血压、血糖、体温、呼吸频率、活动量和睡眠状态等信息，再经过设备端处理、数据传输和云端汇总，形成可用于后续分析的数据。不同参数反映的健康含义并不同，心率和心电主要用于观察心血管状态，血氧和呼吸频率与呼吸功能相关，血糖和血压常用于慢性病管理，活动量和睡眠状态则反映用户的日常生活和恢复情况。

如图10.6所示，可穿戴设备的生理参数采集与传输通常包括传感采集、信号预处理、本地处理、数据传输和云端汇总等环节。传感采集负责从人体获取原始信号，例如光电容积脉搏波、心电信号、连续血糖数据和惯性测量数据；信号预处理用于减小噪声干扰，并完成基础格式整理；本地处理可以在设备端进行初步计算和临时存储；数据传输负责将监测结果同步到手机、网关或云端平台；云端汇总则为后续健康状态识别、异常发现和长期趋势分析提供数据基础。

\begin{figure}[H]
\centering
\begin{tcolorbox}[width=0.88\textwidth, colback=gray!10, colframe=gray!50, boxrule=0.5pt, arc=2pt]
\centering
\vspace{30pt}
{\large\color{gray!70} [图片占位]}
\vspace{10pt}
{\small\color{gray!50} 图10.6 可穿戴设备生理参数采集与传输流程示意图}
\vspace{25pt}
\end{tcolorbox}
\caption*{\footnotesize\itshape\color{gray!50} 图10.6 可穿戴设备生理参数采集与传输流程示意图}
\end{figure}

\subsubsection*{\textbf{1. 可穿戴设备的基本类型与功能}}

可穿戴设备是指可佩戴或贴附在人体表面，并能够较长时间采集生理或行为信息的电子设备。根据佩戴位置和功能用途的不同，常见可穿戴设备可分为腕戴类、胸贴类、皮下植入类、袖带类和其他专用设备等类型。不同设备面对的监测任务不同，采集的数据类型、佩戴方式和适用人群也存在差异。

腕戴类设备包括智能手表和智能手环，具有体积小、佩戴方便、使用门槛低等特点，常用于采集心率、血氧饱和度、活动量和睡眠状态等信息。如图10.7所示，智能手表可以实时显示用户心率，并在后台形成心率、步数和睡眠等健康数据记录。这类设备适合日常健康管理，也常作为连续监测的基础入口。

\begin{figure}[H]
\centering
\begin{tcolorbox}[width=0.88\textwidth, colback=gray!10, colframe=gray!50, boxrule=0.5pt, arc=2pt]
\centering
\vspace{30pt}
{\large\color{gray!70} [图片占位]}
\vspace{10pt}
{\small\color{gray!50} 图10.7 使用智能手表实时监测心率}
\vspace{25pt}
\end{tcolorbox}
\caption*{\footnotesize\itshape\color{gray!50} 图10.7 使用智能手表实时监测心率}
\end{figure}

胸贴类设备主要用于心电信号的长期记录。与腕戴设备相比，胸贴式心电设备更接近临床心电监测场景，能够连续采集单导联或多导联心电信号，适合捕捉阵发性心律失常等短暂出现的异常事件。如图10.8所示，心电贴片通常贴附在胸前，通过内置电极记录心脏电活动，并将心电波形保存或上传至分析系统，为心律异常识别提供数据基础。

\begin{figure}[H]
\centering
\begin{tcolorbox}[width=0.88\textwidth, colback=gray!10, colframe=gray!50, boxrule=0.5pt, arc=2pt]
\centering
\vspace{30pt}
{\large\color{gray!70} [图片占位]}
\vspace{10pt}
{\small\color{gray!50} 图10.8 胸贴式心电监测设备示意图}
\vspace{25pt}
\end{tcolorbox}
\caption*{\footnotesize\itshape\color{gray!50} 图10.8 胸贴式心电监测设备示意图}
\end{figure}

皮下植入类设备的代表是连续血糖监测仪。它通过贴附在上臂或腹部皮肤表面的传感器，将微型探头置入皮下组织间液中，连续检测葡萄糖浓度变化。如图10.9所示，连续血糖监测仪不仅可以显示当前血糖值，还能通过手机端呈现血糖变化曲线，使糖尿病患者的血糖管理从间断的指尖采血转向连续动态观察。

\begin{figure}[H]
\centering
\begin{tcolorbox}[width=0.88\textwidth, colback=gray!10, colframe=gray!50, boxrule=0.5pt, arc=2pt]
\centering
\vspace{30pt}
{\large\color{gray!70} [图片占位]}
\vspace{10pt}
{\small\color{gray!50} 图10.9 连续血糖监测仪及血糖趋势显示示意图}
\vspace{25pt}
\end{tcolorbox}
\caption*{\footnotesize\itshape\color{gray!50} 图10.9 连续血糖监测仪及血糖趋势显示示意图}
\end{figure}

袖带类设备主要用于动态血压监测。设备通过定时充气测量获得24小时血压变化曲线，可用于评价血压控制效果，发现夜间血压异常、晨峰血压和降压药物作用不稳定等问题。如图10.10所示，动态血压监测设备通常由上臂袖带和便携式记录器组成，患者可在日常活动和睡眠过程中完成连续血压记录，从而获得比门诊单次测量更完整的血压变化信息。

\begin{figure}[H]
\centering
\begin{tcolorbox}[width=0.88\textwidth, colback=gray!10, colframe=gray!50, boxrule=0.5pt, arc=2pt]
\centering
\vspace{30pt}
{\large\color{gray!70} [图片占位]}
\vspace{10pt}
{\small\color{gray!50} 图10.10 袖带式动态血压监测设备示意图}
\vspace{25pt}
\end{tcolorbox}
\caption*{\footnotesize\itshape\color{gray!50} 图10.10 袖带式动态血压监测设备示意图}
\end{figure}

除上述设备外，还有一些面向特定场景的专用监测设备。如图10.11所示，睡眠监测设备可用于记录睡眠过程中的呼吸、血氧和睡眠阶段变化，辅助睡眠呼吸暂停等问题的筛查；跌倒监测器主要用于老年人安全管理，在跌倒发生时及时发出提醒；血氧监测指环可连续记录血氧饱和度和脉率，适用于慢性肺部疾病、心衰等人群；呼吸监测腰带则可记录胸腹部运动变化，用于观察呼吸频率和呼吸模式。

\begin{figure}[H]
\centering
\begin{tcolorbox}[width=0.88\textwidth, colback=gray!10, colframe=gray!50, boxrule=0.5pt, arc=2pt]
\centering
\vspace{30pt}
{\large\color{gray!70} [图片占位]}
\vspace{10pt}
{\small\color{gray!50} 图10.11 专用可穿戴监测设备类型示意图}
\vspace{25pt}
\end{tcolorbox}
\caption*{\footnotesize\itshape\color{gray!50} 图10.11 专用可穿戴监测设备类型示意图}
\end{figure}

在功能上，可穿戴设备通常具备生理参数采集、本地数据存储、无线数据传输和简要状态显示等基本能力。部分设备还集成了简单的异常提示、健身指导和用药提醒等附加功能。设备的电池续航能力、数据存储容量、抗干扰水平和佩戴舒适度，是影响其在长期监测场景中实际表现的重要因素。

\subsubsection*{\textbf{2. 生理信号采集的基本原理}}

可穿戴设备采集生理信号所依赖的传感原理，与所测参数的物理特性密切相关。心率监测在腕戴类设备中通常采用光电容积脉搏波技术，即通过发光二极管向皮肤发射特定波长的光，由光电探测器接收被组织反射或透射回来的光信号，根据血液在心脏搏动过程中对光吸收的周期性变化推算心率。该方法具有无创、便捷和适合长期佩戴等优点，但容易受到运动、肤色和佩戴松紧度等因素影响，需要结合滤波和算法处理提高准确性。

心电信号的采集主要依赖生物电极。心脏在每次搏动过程中会在体表产生微弱的电位变化，通过位于胸前或腕部的电极对相应的电位差进行测量，可以记录心电波形。胸贴式心电监测设备通常采用一至多个电极组合，能够长期捕捉心电信号变化，对阵发性心律失常具有较好的发现能力。腕表式心电功能则通常采用单导联设计，可在用户主动触发的方式下完成短时心电记录。

血氧饱和度监测多采用与心率监测类似的光电原理，根据血液中氧合血红蛋白与还原血红蛋白对不同波长光线吸收能力的差异，估算血液的氧合水平。连续血糖监测主要依赖葡萄糖氧化酶电化学传感器，将组织间液中的葡萄糖浓度转化为电流信号。动态血压监测则采用气囊袖带，通过定时充放气过程结合压力传感器记录血压变化。运动状态、体位姿态和跌倒事件等行为信息，通常依赖加速度计、陀螺仪和地磁传感器等惯性测量单元采集相关数据。

\subsubsection*{\textbf{3. 连续监测的特点}}

与院内单次或定时测量相比，连续监测则能提供较长时间范围内的变化曲线。许多临床问题更依赖趋势信息和动态特征，而不是某一次的数值。连续监测不仅关注某个数值是否正常，还关注这个数值什么时候升高、持续多久、是否反复出现，以及是否与睡眠、运动、进食或用药有关。阵发性房颤可能在常规心电图中没有被记录，却能在长期心电监测中被捕捉；夜间血压异常可能在白天测量时并不明显，却能在24小时动态血压曲线中显示出来；糖尿病患者的血糖管理，也不再只依赖几个指尖采血结果，而是可以观察全天血糖的升降变化。

这种连续观察使健康评估从“看一个点”转向“看一条线”。对于慢性心衰患者，心率、活动量和睡眠变化可以反映身体负荷和恢复情况；对于慢性阻塞性肺疾病患者，血氧和呼吸频率的持续变化有助于发现病情加重的早期信号；对于老年高血压患者，不同时段的血压变化能够为用药时间和剂量调整提供参考。

然而，连续监测也会带来新的问题。设备每天都会产生大量数据，心率曲线、心电波形、血糖趋势和活动记录如果全部依靠人工查看，效率很低，也容易遗漏关键信息。因此，后续还需要借助人工智能方法对这些数据进行整理、识别和解释，从连续变化中找出有价值的健康信息。

\subsubsection*{\textbf{4. 数据传输与同步机制}}

可穿戴设备的数据传输通常包括几个环节。设备端负责采集和初步处理数据，并将结果暂存在本地；个人终端负责接收设备上传的信息，如智能手机、平板或专用网关；云端平台负责集中存储、长期管理和进一步分析；在用户授权的前提下，相关结果还可以同步给医疗机构或签约医生，用于健康评估和随访管理。

在这一过程中，数据连续性和安全性十分重要。设备电量不足、蓝牙连接中断、网络不稳定或佩戴脱落，都可能造成数据暂时缺失。系统通常需要通过本地缓存、断点续传和数据校验等方式，减少数据丢失和重复上传。同时，心率、血压、血糖、睡眠等数据属于个人健康信息，传输和存储过程应进行身份认证、权限管理和加密保护，避免数据被误用或泄露。

多设备联合监测时，还需要注意时间同步。心电、血氧、活动量和睡眠数据如果来自不同设备，必须按照统一时间基准对齐。只有这样，系统才能判断某一次心率升高是否发生在运动之后，某一次血氧下降是否出现在睡眠过程中，某一次血糖波动是否与进食或用药有关。时间戳、设备校时和云端同步，是保证多参数联合分析可靠的重要基础。

\subsubsection*{\textbf{5. 小结}}

本节介绍了生理参数采集与连续监测的基本过程。智能手表、心电贴片、连续血糖监测仪、动态血压监测设备和各类专用监测设备，能够在日常生活中持续采集心率、血氧、心电、血糖、血压、活动和睡眠等信息。这些数据经过设备端、个人终端、云端平台和医疗系统的传输与同步后，形成可用于后续分析的连续健康记录，为健康状态识别、异常发现和风险预警提供基础。

\subsection{健康状态识别与异常发现}

心率曲线、心电波形、血氧变化、血糖曲线、血压趋势、活动量和睡眠记录，看起来都是一条条不断变化的线。对普通用户来说，它们可能只是屏幕上的高低起伏；对健康评估系统来说，这些曲线中包含着身体状态变化的线索。系统需要从这些连续数据中提取特征，判断当前状态是否稳定，并发现可能需要关注的异常变化。

如图10.12所示，健康状态识别与异常发现通常包括特征提取、状态识别、异常识别、个体化基线分析和多参数融合等环节。系统首先从心率、心电、血氧、血糖、血压、活动量和睡眠等数据中提取时域特征、频域特征以及形态特征；随后结合阈值规则、传统机器学习和深度学习方法，对健康状态进行分级判断，并识别异常波动；最后通过多参数融合输出风险等级、趋势曲线和异常时间窗，为用户提醒和医生复核提供依据。

\begin{figure}[H]
\centering
\begin{tcolorbox}[width=0.88\textwidth, colback=gray!10, colframe=gray!50, boxrule=0.5pt, arc=2pt]
\centering
\vspace{30pt}
{\large\color{gray!70} [图片占位]}
\vspace{10pt}
{\small\color{gray!50} 图10.12 健康状态识别与异常发现流程示意图}
\vspace{25pt}
\end{tcolorbox}
\caption*{\footnotesize\itshape\color{gray!50} 图10.12 健康状态识别与异常发现流程示意图}
\end{figure}

\subsubsection*{\textbf{1. 健康状态识别的基本任务}}

健康状态识别是指系统根据连续监测数据，对用户当前或近期的身体状态进行判断。常见任务包括心律状态识别、血压控制水平评估、血糖波动分析、睡眠状态判断、活动强度识别和跌倒风险判断等。不同任务关注的数据不同，但目标都是把复杂的生理信号转化为更容易理解的健康信息。

在实际应用中，健康状态识别通常包括两类结果。一类是状态判断，如正常心律或疑似心律异常、清醒或睡眠、静息或运动、低风险或高风险等；另一类是指标评估，如平均心率、血压负荷、血糖达标时间、睡眠效率和活动水平等。这些结果可以帮助用户了解自身状态，也能为医生随访和健康管理提供参考。

需要注意的是，健康状态识别不能只看单个数值。运动后心率升高可能是正常反应，静息状态下心率持续偏高才更值得关注；血糖升高需要结合进餐时间和用药情况判断；睡眠中血氧下降也要结合持续时间和发生频率分析。

\subsubsection*{\textbf{2. 特征提取与状态判断}}

可穿戴设备采集到的数据往往表现为连续曲线或波形。系统不能只盯着某一个数字，而要从这些曲线中提取有意义的特征。所谓特征，可以理解为能够描述健康状态的关键信息，例如心率的平均值、波动幅度、变化速度，血糖曲线的峰值和持续时间，睡眠过程中的觉醒次数和深睡比例等。

特征提取通常包括时域特征、频域特征和形态特征。时域特征直接描述指标在时间上的变化，如最大值、最小值、均值和持续时间；频域特征用于分析信号的节律变化，常见于心率变异性和睡眠呼吸分析；形态特征则关注波形形状，如心电波形中的QRS波群、ST段变化和异常节律等。

在人工智能系统中，传统机器学习方法可以利用这些人工设计的特征完成分类和预测，深度学习方法则可以从原始波形或时间序列中自动学习复杂模式。例如，卷积神经网络可以识别局部波形特征，循环神经网络和长短期记忆网络适合处理时间变化，Transformer结构可以捕捉较长时间范围内的关联。不同方法各有特点，实际应用中常根据数据类型和任务目标进行选择。

\subsubsection*{\textbf{3. 异常识别方法}}

异常识别是指系统从连续监测数据中发现偏离正常状态或可能提示健康风险的变化。可穿戴设备采集的数据通常具有时间连续、波动频繁和个体差异明显等特点，因此异常识别不能只依赖某一个瞬时数值，而需要结合指标范围、持续时间、变化趋势、发生场景和个体基线进行综合判断。

阈值与规则判断是最基础的异常识别方法。系统可以根据医学参考范围、疾病管理目标或专家规则设定判断条件。例如，血氧饱和度低于设定范围、血糖持续高于或低于目标区间、血压在多个时间点持续偏高，均可触发异常提示。这类方法逻辑清楚、解释性强，适合处理边界明确、医学规则相对稳定的指标。但它对个体差异的适应能力有限，容易将运动、情绪波动或短暂测量干扰误判为异常。

统计分析方法主要根据数据分布和变化规律识别异常。系统可以计算某一指标的均值、标准差、波动范围和变化趋势，并判断当前数据是否明显偏离历史水平。例如，某位用户平时静息心率较稳定，如果连续数日静息心率明显升高，即使没有超过通用阈值，也可能提示身体状态发生变化。统计方法适合发现趋势性变化和个体偏离，但需要一定数量的历史数据作为参考。

机器学习方法可以利用多种特征进行异常识别。系统先从原始数据中提取心率变化、心率变异性、血氧下降次数、血糖波动幅度、睡眠时长、活动量变化等特征，再利用分类模型或异常检测模型判断状态是否异常。常见方法包括支持向量机、随机森林、梯度提升树、聚类分析和孤立森林等。这类方法能够处理多个指标之间的组合关系，适合用于心律异常、睡眠异常、血糖波动异常和活动模式异常等任务。

深度学习方法适合处理长时间序列和复杂波形数据。对于心电信号、连续血糖曲线、呼吸信号和多日活动数据，深度学习模型可以直接从连续数据中学习变化模式。卷积神经网络适合提取局部波形特征，循环神经网络和长短期记忆网络适合分析前后时间依赖关系，Transformer结构则适合捕捉较长时间范围内的变化关联。这类方法在复杂数据分析中具有优势，但通常需要较多训练数据，并且结果解释需要结合可视化和医学规则进一步说明。

不同类型的异常具有不同特点，识别方法也需要相互配合。血氧过低、血糖明显异常和血压持续偏高等情况，通常适合先用阈值规则进行初步判断；心律异常、睡眠呼吸异常和复杂血糖波动等问题，则需要结合机器学习或深度学习模型进行进一步分析；长期健康管理还需要结合个体化基线和趋势变化，判断某一指标偏离是否具有持续风险。通过多种方法协同，系统既能保留规则判断的可解释性，又能提高复杂监测场景下的异常识别能力。

\subsubsection*{\textbf{4. 个体化基线与动态评估}}

同一个数值，对不同人可能有不同意义。有些人平时静息心率较低，有些人血压昼夜波动明显，有些糖尿病患者餐后血糖变化幅度较大。例如，运动员的静息心率常显著低于普通人，心率50次/分钟在一位马拉松爱好者身上是正常状态，在一位70岁的老人身上则可能是值得警惕的信号。如果系统完全按照统一阈值判断，就容易把正常个体差异误认为异常，也可能忽略某些相对变化带来的风险。

个体化基线是指根据用户长期监测数据建立个人参考范围。系统可以记录用户在不同时间段、不同活动状态和不同睡眠状态下的心率、血压、血糖、血氧和活动水平，形成相对稳定的个人健康轨迹。后续监测中，系统不仅判断指标是否超过通用范围，还会观察它是否明显偏离用户自身的长期状态。

动态评估强调趋势变化，而不是只看某一刻的结果。一次心率升高可能只是运动后的正常反应，但连续几天静息心率升高、睡眠变差和活动量下降同时出现，就可能提示健康状态发生变化。通过个体化基线和趋势分析，系统可以更早发现风险线索，使健康评估更加贴近日常生活中的真实状态。

\subsubsection*{\textbf{5. 多参数融合与结果解释}}

人体状态通常不会只反映在一个指标上。心率升高可能与运动有关，也可能与发热、焦虑或心律异常有关；血氧下降可能是呼吸问题，也可能是佩戴不稳造成的测量偏差；睡眠质量下降可能与疼痛、压力、呼吸异常或生活习惯有关。单一参数很容易产生误判，多参数融合可以帮助系统更全面地理解健康状态。

多参数融合是指将心率、心电、血氧、血压、血糖、活动、睡眠和体温等信息结合起来分析。例如，心率升高同时伴随步数增加，通常可以解释为运动反应；如果静息状态下心率升高并伴随血氧下降，就需要进一步关注；如果连续几天活动量下降、睡眠变差并伴随心率变化，也可能提示身体状态出现波动。

结果解释同样重要。系统不能只给出“异常”两个字，而应说明异常发生在什么时间、涉及哪些指标、持续多久、与个人平时状态相比变化多大。清楚的解释帮助用户理解提醒内容，也方便医生判断是否需要进一步检查、随访或干预。

\subsubsection*{\textbf{6. 小结}}

本节介绍了健康状态识别与异常发现的基本过程。可穿戴设备采集到的心率、心电、血氧、血压、血糖、活动量和睡眠等连续数据，需要经过特征提取、状态判断、异常识别和个体化分析，才能转化为可理解的健康信息。阈值规则适合识别边界清楚的异常，统计方法可以发现个体长期变化趋势，机器学习和深度学习方法则能够处理更复杂的波形和多参数关系。通过这些方法的配合，人工智能可以从连续监测数据中发现潜在健康风险，并为用户提醒、医生复核和后续健康管理提供依据。

\subsection{案例：基于可穿戴设备的健康状态监测}

慢性心衰患者出院以后，真正的管理并没有结束。白天活动量下降、夜间心率升高、体重短时间增加、血氧反复下降，这些变化可能出现在家里，也可能发生在睡眠中。患者自己未必能及时察觉，医生也不可能每天面对面观察。可穿戴设备和远程健康管理系统的价值，正在于把这些院外变化连续记录下来，并在风险积累到一定程度时及时提醒医护人员和患者。

本案例以一名72岁慢性心衰患者的90天远程健康管理为例。患者为女性，确诊慢性心力衰竭多年，近期曾因急性失代偿入院治疗，出院后被纳入医院心衰中心远程健康管理项目。系统通过智能手表、单导联心电贴片、动态血压监测设备、家用血压计和体重秤等设备，连续采集患者心率、活动量、心电片段、血压、血氧和体重等数据，并将其上传至云端平台。心衰中心护士和签约医师根据平台生成的预警信息和健康报告，对患者进行随访、用药指导和就医建议。

如图10.13所示，本案例中的健康状态监测系统由用户端、云端平台和医护端组成。用户端负责采集居家场景中的连续健康数据，云端平台负责完成数据清洗、健康档案构建、模型分析和分级预警，医护端则根据系统生成的健康报告进行复核、随访和干预。

\begin{figure}[H]
\centering
\begin{tcolorbox}[width=0.88\textwidth, colback=gray!10, colframe=gray!50, boxrule=0.5pt, arc=2pt]
\centering
\vspace{30pt}
{\large\color{gray!70} [图片占位]}
\vspace{10pt}
{\small\color{gray!50} 图10.13 基于可穿戴设备的健康状态监测系统总体架构示意图}
\vspace{25pt}
\end{tcolorbox}
\caption*{\footnotesize\itshape\color{gray!50} 图10.13 基于可穿戴设备的健康状态监测系统总体架构示意图}
\end{figure}

监测对象与系统配置

本案例中的患者属于慢性心衰高风险人群。此类患者出院后仍可能出现液体潴留、心律异常、血压波动、血氧下降和活动耐量下降等问题。如果这些变化没有被及时发现，可能导致病情加重甚至再次入院。因此，本案例的监测重点不是单独观察某一个指标，而是连续追踪心衰相关的多项生理和行为变化。

系统为患者配置了多种监测设备。智能手表用于持续记录心率和活动量，单导联心电贴片用于每日固定时间段采集短时心电波形，动态血压监测设备用于每周两次完成24小时血压记录，家用血压计和体重秤用于补充日常血压和体重测量。所有数据通过家庭网关上传至云端平台，并与患者电子健康档案关联。

云端平台负责对数据进行汇总和分析，形成连续健康档案。心衰中心护士主要负责日常数据查看和初步随访，签约医师负责对中高风险预警进行医学判断，并根据患者情况调整治疗方案。这样，患者居家监测、云端分析和医护复核之间形成了较完整的远程管理链条。

2. 90天远程监测过程

在90天观察期内，系统持续记录患者的心率、活动量、血压、心电片段和体重变化。与传统复诊相比，这种监测方式能够覆盖患者出院后的日常生活状态。系统不仅能看到某一天的血压或心率是否异常，还能观察这些指标在几天或几周内是否持续偏离个人基线。

监测期间，平台累计记录心率数据约130万条、心电片段约600段、24小时动态血压记录24次、家用血压测量约180次和体重测量约90次，并生成每周健康报告13份。这些数据使医护人员能够较连续地了解患者出院后的健康变化，而不是只依赖患者复诊时的主观描述和少量检查结果。

系统在数据分析时重点关注三类变化。第一，体重和心率的共同变化，用于提示液体潴留和心衰加重风险。第二，心电片段和活动量之间的关系，用于判断心律异常是否与运动有关。第三，夜间血氧和呼吸频率变化，用于发现睡眠期间可能出现的呼吸或循环风险。通过这些分析，系统在90天内共触发3次具有临床意义的预警事件。

第一次预警：液体潴留与心衰加重风险

第一次预警发生在出院第18天。系统发现患者夜间出现持续性快速心率，同时体重在3天内增加约2kg。对于慢性心衰患者而言，短时间体重增加常提示体液潴留风险；如果同时伴随心率升高，就需要警惕心衰加重的早期表现。

平台将该事件判定为中度预警，并将信息同步给签约医师。医师结合患者电子健康档案、近期体重变化、心率趋势和用药记录后，认为患者可能存在液体潴留和心衰加重的早期倾向，于是在原有治疗基础上调整利尿剂剂量，并安排患者门诊复诊。

患者按建议执行后，体重逐步回落，夜间心率也逐渐恢复至个人基线范围。该次预警说明，连续监测可以在患者症状明显加重之前捕捉到风险变化，为医生提前调整治疗方案提供依据。

第二次预警：阵发性心律异常识别

第二次预警发生在出院第45天。系统在每日心电片段中识别出多次阵发性心律异常，并结合活动量数据进行判断。由于异常发生时患者并未处于明显活动状态，系统认为该异常不太可能单纯由运动引起，因此触发中度预警。

签约医师通过远程会诊查看心电波形和近期趋势数据后，建议患者到医院接受进一步检查。后续检查确认患者存在阵发性房颤。根据这一结果，医师对其治疗方案进行了调整，并补充了抗凝治疗相关管理。

这一事件体现了可穿戴心电监测的优势。阵发性心律异常可能在常规门诊心电图中没有出现，却能在居家连续监测过程中被多次捕捉。系统通过心电片段、活动量和时间信息综合判断，帮助医生更快定位需要复核的异常线索。

第三次预警：夜间血氧下降与及时就医

第三次预警发生在出院第72天。系统监测到患者夜间血氧饱和度明显下降，并伴有呼吸频率升高。由于这一变化发生在睡眠期间，患者本人可能并不能及时察觉。系统将该事件判定为重度告警，并同时通知患者、家属和签约医师。

家属收到通知后，陪同患者前往附近医院急诊就诊。经评估后，患者接受住院观察和治疗。该次预警使患者在病情进一步恶化前获得及时处理，避免了一次潜在的失代偿性心衰再入院事件。

与前两次中度预警不同，第三次事件体现了分级告警机制的重要性。对于可能涉及急性风险的异常，系统不能只在应用程序中给出普通提醒，而需要通过多渠道通知患者、家属和医护人员，并引导其尽快进入医疗处置流程。

\subsubsection*{\textbf{3. 管理结果与应用价值}}

在整个90天远程管理过程中，系统共触发3次具有临床意义的预警。其中，2次直接促成治疗方案调整，1次促成及时就医并避免潜在再入院事件。患者在管理结束后对系统使用体验进行了评价，对设备佩戴舒适度、提示频率和医护响应及时性给予较高反馈，并表示愿意继续参与后续远程管理。

从结果看，本案例的价值主要体现在三个方面。第一，系统能够把出院后的零散健康变化转化为连续记录，使医护人员及时发现体重增加、心率异常、心律异常和血氧下降等风险信号。第二，系统能够将异常事件与患者个人基线、活动状态和既往病史联系起来，减少单一指标判断带来的误差。第三，系统能够把预警结果推送给不同角色，使患者、家属、护士和签约医师围绕同一份健康报告进行协同处理。

该案例说明，可穿戴设备的作用不只是采集数据，而是参与到慢性病院外管理的完整流程中。数据采集、云端分析、分级预警、医护复核和干预反馈共同构成了远程健康管理闭环。对于慢性心衰这类需要长期随访的疾病，可穿戴监测可以帮助医生更早发现病情波动，也能帮助患者在日常生活中获得更及时的健康支持。

\subsubsection*{\textbf{4. 小结}}

本案例以慢性心衰患者90天远程健康管理为例，展示了可穿戴设备在健康状态监测中的应用过程。系统通过智能手表、心电贴片、动态血压监测设备和家用监测设备持续采集心率、活动量、心电、血压、血氧和体重等数据，并利用云端平台进行趋势分析和分级预警。在观察期内，系统共记录多源健康数据并生成13份周报告，触发3次具有临床意义的预警，其中2次促成治疗方案调整，1次促成及时就医并避免潜在再入院事件。该案例表明，智能穿戴监测能够把院外连续数据转化为可识别、可复核和可干预的健康管理信息，为慢性病远程随访和个体化管理提供支持。

\section{智能康复系统与功能训练}

康复训练室里，一名脑卒中后患者正在练习抬手。过去，治疗师主要依靠观察和经验判断患者的动作是否到位：手臂抬得够不够高，手指能不能张开，步态是否对称，训练后有没有疲劳。然而，康复训练不是做一次检查就能完成的事情，它往往需要持续数周甚至数月，每一次动作的幅度、力量、速度和稳定性，都可能反映患者功能恢复的变化。仅靠人工观察，很难完整记录这些细节。

智能康复系统的作用，是把康复训练过程变得更加可观察、可记录和可调整。系统通过康复机器人、外骨骼、智能手套、平衡训练设备、传感器和康复管理平台等工具，采集患者训练过程中的关节角度、肌肉用力、步态变化、动作完成度和训练疲劳情况。人工智能方法则进一步分析这些数据，帮助治疗师评估功能状态、设计训练方案、调整训练强度，并对康复效果进行阶段性评价。

如图10.14所示，智能康复系统与功能训练通常包括康复设备硬件、人机交互控制、训练方案设计、过程数据采集和效果评估反馈等。康复设备为患者提供训练支持，人机交互控制保障训练过程安全、灵活和可调，训练方案设计用于确定任务、强度和阶段安排，过程数据采集记录训练中的轨迹、力量和完成度，效果评估反馈则用于判断训练效果并指导下一阶段方案调整。整个过程形成“评估—调整—再训练”的循环，使康复训练从经验驱动逐步转向数据支持和人机协同。

\begin{figure}[H]
\centering
\begin{tcolorbox}[width=0.88\textwidth, colback=gray!10, colframe=gray!50, boxrule=0.5pt, arc=2pt]
\centering
\vspace{30pt}
{\large\color{gray!70} [图片占位]}
\vspace{10pt}
{\small\color{gray!50} 图10.14 智能康复系统与功能训练总体框架示意图}
\vspace{25pt}
\end{tcolorbox}
\caption*{\footnotesize\itshape\color{gray!50} 图10.14 智能康复系统与功能训练总体框架示意图}
\end{figure}

\subsection{康复设备构成与交互控制}

康复设备及其交互控制是智能康复系统能够安全运行的基础。设备硬件负责感知患者状态和执行训练动作，控制系统负责调节训练方式，人机交互界面负责向患者和治疗师呈现任务与反馈，安全保护机制则保证训练过程不超过患者承受范围。如图10.15所示，康复设备通常由感知模块、执行模块、控制模块、交互模块和安全保护模块组成，同时支持被动训练、主动训练、辅助训练和阻抗训练等多种人机交互方式。设备通过感知患者动作状态，调节辅助力或阻力，并将训练结果反馈给患者和治疗师，从而实现较为灵活的功能训练。

\begin{figure}[H]
\centering
\begin{tcolorbox}[width=0.88\textwidth, colback=gray!10, colframe=gray!50, boxrule=0.5pt, arc=2pt]
\centering
\vspace{30pt}
{\large\color{gray!70} [图片占位]}
\vspace{10pt}
{\small\color{gray!50} 图10.15 康复设备组成与交互控制示意图}
\vspace{25pt}
\end{tcolorbox}
\caption*{\footnotesize\itshape\color{gray!50} 图10.15 康复设备组成与交互控制示意图}
\end{figure}

\subsubsection*{\textbf{1. 感知模块与状态采集}}

感知模块相当于康复设备的“眼睛”和“触觉”。它负责采集患者在训练过程中的动作、力量、姿态和肌肉活动等信息，使系统能够知道患者正在怎样运动。常见感知信息包括关节角度、肢体位置、运动速度、加速度、足底压力、握力、肌电信号和身体重心变化等。

在上肢康复训练中，感知模块可以记录肩、肘、腕等关节的活动范围和运动轨迹，帮助判断患者是否完成抬臂、伸手、抓握等动作。在下肢康复训练中，系统可以通过压力传感器和惯性传感器记录步态、支撑时间、步幅和重心转移情况。在肌肉功能训练中，表面肌电传感器可以反映肌肉是否正确激活，帮助识别患者是否存在发力不足或异常协同。

这些数据不仅用于记录训练结果，也用于实时控制。如果系统检测到患者发力不足，可以增加辅助；如果发现动作偏离目标轨迹，可以及时提醒；如果出现异常角度或用力过大，则需要触发保护机制。

\subsubsection*{\textbf{2. 执行模块与训练驱动}}

执行模块负责产生训练所需的运动或力量，是康复设备真正“动起来”的部分。常见执行方式包括电机驱动、气压驱动、液压驱动和弹性辅助结构等。不同执行方式适用于不同训练场景。电机驱动控制精度较高，常用于上肢机器人、手部训练设备和步态训练系统；气压驱动柔顺性较好，适合需要轻柔辅助的训练；液压驱动输出力量较大，多用于承重或下肢训练设备。

在功能训练中，执行模块可以提供助力，也可以提供阻力。对于早期康复患者，设备可以带动肢体完成关节活动，帮助维持活动范围，减少僵硬和挛缩风险。对于恢复期患者，设备可以根据患者主动发力情况提供适当辅助，使其完成目标动作。对于功能强化阶段，设备还可以设置阻力，帮助患者增强肌力和控制能力。

执行模块的设计必须重视平稳性和安全性。康复训练面对的是功能受损患者，而不是健身房里兴奋过头的人类。设备输出的速度、力量和活动范围都需要受到严格限制，并根据患者状态动态调整，避免造成牵拉、疼痛或二次损伤。

\subsubsection*{\textbf{3. 控制模块与训练调节}}

控制模块是康复设备的“大脑”，负责把感知模块采集到的信息转化为具体控制指令。它需要根据患者当前动作、训练目标和安全边界，调节执行模块的运动方向、速度、力量和辅助程度。控制模块的好坏，直接决定设备是“顺着患者训练”，还是“强行带着患者运动”。

常见控制方式包括位置控制、力控制、阻抗控制和自适应控制。位置控制主要保证肢体按照设定轨迹运动，适用于被动训练或早期关节活动训练。力控制关注设备输出的力量大小，适用于力量训练或辅助训练。阻抗控制则在位置和力量之间保持柔顺关系，使设备既能提供支持，又不会过度限制患者主动动作。自适应控制可以根据患者训练表现动态调整参数，更适合个体化康复训练。

在实际训练中，控制模块需要不断根据患者状态调整训练方式。例如，患者主动发力增强时，设备可以减少辅助，让患者承担更多控制；患者出现疲劳或动作完成质量下降时，设备可以降低难度或增加辅助；患者动作偏离安全范围时，设备应立即减速或停止。这样的调节能力，是智能康复设备区别于普通机械训练器械的重要特征。

\subsubsection*{\textbf{4. 交互模块与训练反馈}}

交互模块负责让患者和治疗师看得懂、用得上。康复训练不能只是设备内部运算，患者需要知道自己要做什么，治疗师也需要知道患者完成得怎么样。因此，交互模块通常包括显示屏、操作面板、语音提示、虚拟训练场景、触觉反馈和训练报告等内容。

对于患者来说，交互模块可以把枯燥的重复训练转化为可理解的任务。例如，屏幕上显示目标位置，患者通过抬手、伸手或移动重心完成任务；系统用得分、进度条或轨迹图提示动作完成情况；语音反馈提醒患者调整姿势或控制节奏。虚拟现实训练还可以把动作任务设计成游戏化场景，提高训练参与度。

对于治疗师来说，交互模块提供的是训练管理信息。系统可以显示患者的动作轨迹、训练次数、完成率、关节活动范围、肌力变化和疲劳趋势。治疗师据此判断训练是否合适，是否需要调整任务难度、辅助力或训练时间。

\subsubsection*{\textbf{5. 安全保护与风险控制}}

安全保护模块是康复设备中最不能省略的部分。康复患者可能存在肌力不足、关节稳定性差、疼痛敏感、平衡能力下降等问题，训练过程中如果设备控制不当，可能带来牵拉、跌倒、疼痛或疲劳加重等风险。因此，智能康复系统必须在硬件和软件两个层面设置安全边界。

硬件保护通常包括急停按钮、机械限位、扭矩限制、角度限制、支撑结构和防跌倒装置等。软件保护则包括异常检测、速度限制、力量阈值、训练中断策略和报警提示等。当系统检测到患者动作异常、力量过大、关节角度超限或设备运行故障时，应及时停止训练或切换到安全模式。

安全控制还需要结合患者个体情况。不同患者的关节活动范围、肌力水平和耐受程度不同，安全阈值也不应完全一致。系统应根据治疗师设定和患者训练表现动态调整保护参数。康复训练的目标是促进功能恢复，而不是把患者训练成勇闯机械迷宫的倒霉角色。

\subsubsection*{\textbf{6. 人机交互训练方式}}

智能康复设备常见的人机交互训练方式包括被动训练、主动训练、辅助训练和阻抗训练。被动训练主要用于康复早期，患者主动运动能力较弱，设备带动肢体完成关节活动，帮助维持活动范围。主动训练适用于患者具备一定运动能力时，设备主要记录动作并提供反馈，鼓励患者自主完成训练任务。

辅助训练介于被动训练和主动训练之间。当患者能够主动发力但力量不足时，设备在关键阶段提供助力，帮助其完成完整动作。阻抗训练多用于恢复期或强化期，设备提供一定阻力，促进肌力、耐力和协调控制能力提升。四种方式并不是互相割裂的，而是可以随着患者恢复进展逐步切换。

合理的人机交互控制应遵循“按需辅助”的原则。辅助过多，患者容易依赖设备，主动参与不足；辅助过少，患者可能无法完成动作，甚至产生挫败感或代偿动作。智能康复系统需要根据训练表现动态调节辅助程度，使患者始终处于安全、可完成且具有挑战性的训练状态。

\subsubsection*{\textbf{7. 小结}}

本节介绍了智能康复设备的基本构成与交互控制方式。康复设备通常由感知模块、执行模块、控制模块、交互模块和安全保护模块组成，能够采集患者动作状态，提供训练助力或阻力，并通过界面和反馈帮助患者完成训练任务。被动训练、主动训练、辅助训练和阻抗训练分别适用于不同康复阶段，系统通过按需辅助和动态调节，使训练过程更加安全、灵活和个体化。智能康复设备的核心价值不在于让机器替代治疗师，而在于为治疗师提供更连续、更量化和更可反馈的训练支持。

\subsection{康复训练设计与反馈调节}

设备只是舞台，真正决定康复效果的是演什么、怎么演，训练任务的合理设计与训练过程中的反馈调节。同样一个抬手的动作，机械重复100次和每天根据患者状态调整难度、配合即时反馈练 100 次，长期效果可能完全不同。在智能康复系统中，训练设计与反馈调节通常由康复治疗师、医师与人工智能算法共同完成，体现着医工融合在临床应用中的实际价值。

如图10.11所示，康复训练设计与反馈调节通常包括训练目标设定、任务设计、反馈机制构建、过程评估和方案调整等环节，多个环节通过数据闭环连接，形成持续优化的训练过程。

\begin{figure}[H]
\centering
\begin{tcolorbox}[width=0.88\textwidth, colback=gray!10, colframe=gray!50, boxrule=0.5pt, arc=2pt]
\centering
\vspace{30pt}
{\large\color{gray!70} [图片占位]}
\vspace{10pt}
{\small\color{gray!50} 图10.11 康复训练设计与反馈调节闭环示意图}
\vspace{25pt}
\end{tcolorbox}
\caption*{\footnotesize\itshape\color{gray!50} 图10.11 康复训练设计与反馈调节闭环示意图}
\end{figure}

\subsubsection*{\textbf{1. 康复训练的基本原理与设计要素}}

康复训练的科学基础来自两个理论：神经可塑性和运动学习。神经可塑性是指神经系统在外界刺激和反复练习作用下，可以调整自身结构和功能的能力，通俗讲，多走一次，就多通一分路。脑卒中、脑外伤等中枢神经损伤后，部分受损通路虽难以完全恢复，但通过反复、有目的的训练，可以促使未受累区域承担起部分功能，从而改善运动控制和日常生活能力。

运动学习理论进一步指出，技能习得依赖四个条件：明确的任务目标、合理的练习强度、即时的反馈、适当的难度挑战，任意一项缺位，效果都会打折扣。康复训练的科学设计正是建立在这两个理论之上，把重复变成有方向、有节奏、有反馈的反复练习。

在训练方案设计中，通常要考虑三类基本要素。第一是训练目标，比如恢复某一关节活动度、增强某一肌群肌力、改善步态稳定性、训练手部精细操作或提高日常生活能力等，目标应当与患者整体康复计划相一致。

第二是训练强度与频次，包括每次训练时长、重复次数、组间休息和每周训练频次等，决定训练在生理层面的累积效应。第三是训练难度，包括动作复杂程度、施加阻力大小、运动速度要求和任务情境复杂性，决定训练的挑战性。三者关系类似健身房里的动作、组数、负重，目标定方向，强度保积累，难度提挑战。

三类要素之间相互关联。训练目标决定任务的方向，训练强度与频次保证训练效应的累积，训练难度则在合适水平上保持患者的兴趣与挑战感。设计不当容易出现两种极端情形：训练难度过低、强度不足，训练效果难以体现；训练难度过高、强度过大，则可能加重患者疲劳并增加损伤风险。康复训练的科学设计，正是要在患者当前能力与目标功能之间找到合适的平衡。

\subsubsection*{\textbf{2. 训练任务设计与个体化方案}}

智能康复系统的一大优势，在于能根据患者具体情况设计个体化训练方案。个体化方案的形成通常包括基线评估、任务选择、参数设置和阶段安排四个环节。

基线评估在训练开始前对患者的运动功能、肌力、关节活动度、平衡能力和日常生活能力进行综合评价，常用工具包括 Fugl-Meyer评分、Brunnstrom 分期、徒手肌力测试、关节活动度测量、Berg 平衡量表和 Barthel 指数等，这些量表共同回答患者目前处于什么水平、能从哪里起步，是确定训练起点和目标的重要依据。

任务选择是根据基线评估结果确定具体训练内容。对于上肢功能障碍患者，可选择伸够取物、抓握释放、翻书与穿珠等任务；对于下肢与步态训练，可选择平衡转移、跨步、上下楼梯和不同地面行走等任务；对于手部精细动作训练，可选择捏取、按压、对指动作和书写等任务。任务通常以患者实际生活中的动作为参考，使训练更具功能意义。参数设置则在任务选定后确定具体的运动范围、阻力、速度、重复次数和反馈方式等。

在阶段安排方面，康复训练通常按由易到难、由被动到主动、由单关节到多关节、由简单环境到复杂环境的方向逐步推进。智能康复系统可以根据每次训练后的数据，对患者表现进行客观评估，并将评估结果反馈到下一次训练方案中，形成动态调整机制。例如，当某一动作的完成度持续达到目标后，系统可以建议提升任务难度或引入新的训练项目；当患者表现出明显疲劳或动作质量下降时，系统则可以建议降低强度或安排适当休息。

\subsubsection*{\textbf{3. 反馈机制与调节方式}}

反馈是康复训练中不可缺少的组成部分。即时、清晰和恰当的反馈能够帮助患者更准确地理解动作要求、及时纠正动作偏差，并增强训练的参与感。智能康复系统通常通过多种感官通道向患者提供反馈，包括视觉反馈、力反馈、听觉反馈和生理反馈等。

视觉反馈通过显示屏或投影画面将训练动作、轨迹偏差和完成情况以图形方式呈现给患者，常采用游戏化界面或虚拟现实场景增强吸引力。例如，将上肢伸够动作设计为画线、击球或采摘水果等任务，使训练过程具有较强的目标感和趣味性。力反馈通过机械装置向患者肢体提供可感知的阻力或推动力，使其在完成动作时获得真实的力感，常用于阻抗训练和精细动作训练。听觉反馈通过提示音或语音指导提示动作节奏与完成情况，对节律性较强的训练尤其有效。

生理反馈把患者自身的生理信号实时呈现给患者，让他能直接看到和调节自己的身体状态。常见形式包括基于肌电信号的肌肉激活反馈、基于脑电或脑功能信号的神经控制反馈以及基于心率变异性的自主神经调节反馈等，这些原本看不见的生理状态被转化为可感知的信号，帮助患者建立对自身身体活动的更精细控制。

在反馈调节过程中，频率、强度和呈现形式都需要与训练目标和患者认知水平相匹配，反馈过多容易分散注意力，反馈过少又会找不到方向。这个分寸常常需要康复治疗师与系统共同把握，也是智能康复系统区别于自动化器械的关键所在。

\subsubsection*{\textbf{4. 训练过程的智能调节}}

在训练过程中，AI 可以根据连续采集的数据对训练参数进行实时或周期性调节，主要包含四件事：动作质量识别、代偿动作识别、疲劳状态评估和难度自适应调整。动作质量识别依据传感器采集的轨迹、速度和力的信息，对动作的完整性、稳定性和精确度进行评估。

代偿动作识别更有意思，它要发现患者在功能受限时投机取巧的现象，比如本该靠肩部上抬完成的动作，却用躯干前倾代偿；这类代偿如果不及时纠正，长期可能形成不良运动模式，影响功能恢复。

疲劳状态评估依据训练过程中的肌电信号、心率变化、动作节奏和力的稳定性等综合判断患者是否出现疲劳，必要时建议调整训练强度或安排休息。难度自适应调整通过对每次训练数据的分析，将动作难度、阻力大小和任务复杂度与患者当前能力进行匹配，使训练始终保持在“具有挑战性但可以完成”的合适水平。智能调节的目标，是在保证训练安全的前提下，使每一次训练都尽可能贴近患者当前需要。

在算法层面，训练过程的智能调节可结合时间序列分析、模式识别、强化学习和规则推理等多种方法。时间序列分析用于刻画训练过程中各项指标的变化趋势；模式识别用于识别动作类别和异常模式；强化学习用于在多次训练中逐步优化难度调整策略；规则推理用于结合临床经验和安全限制对调整方案进行约束。多种方法结合使用，使智能调节既具有较强的数据适应能力，又能够保留必要的临床合理性。

\subsubsection*{\textbf{5. 数据闭环与远程支持}}

智能康复系统的另一重要价值在于建立训练数据的闭环管理机制。每次训练过程中，设备记录的运动数据、力数据、肌电数据和反馈交互数据等，可以同步上传到康复管理平台，与患者电子健康档案相结合，形成动态、持续的训练记录。康复治疗师和医师可以通过平台查看每位患者的训练进展、参数调整情况和功能评估结果，并据此对训练方案进行修订。

数据闭环通常以一定时间周期为单位组织。例如，以日为单位生成训练日报，记录当日训练任务、参数设置、完成情况和反馈情况；以周为单位生成训练周报，分析训练强度、动作质量变化和功能改善趋势；以月或阶段为单位生成阶段评估报告，结合标准量表评分对康复进展进行较为系统的总结。报告内容通常以图表与文字相结合的形式呈现，便于康复团队进行远程审阅和方案调整。

在远程支持方面，智能康复系统可以扩展到家庭康复和社区康复场景。患者在家中使用便携式康复设备或日常生活中的可穿戴设备进行训练，所采集的数据通过家庭网关或移动网络上传至云端平台，由签约康复治疗师远程查看并提供指导。远程支持模式有助于将康复服务从医院延伸到日常生活中，解决出院后训练频次下降、康复中断和长期效果难以维持等问题，但也对设备稳定性、数据隐私保护和远程沟通机制提出了较高要求。

\subsubsection*{\textbf{6. 小结}}

康复训练设计与反馈调节回答的是三个层次的问题：练什么（训练目标 + 任务设计）、怎么练（强度 + 难度 + 阶段安排）、练得怎样（动作质量 + 代偿识别 + 疲劳评估 + 难度自适应）。视觉、力、听觉、生理四类反馈在训练中边练边纠错，日报、周报、阶段评估的数据闭环让医院、社区和家庭场景下的康复都能被持续追踪。这套体系背后，神经可塑性和运动学习理论提供了科学依据，把机械重复变成有反馈的反复练习，正是智能康复系统的核心价值。

\subsection{案例：康复机器人在脑卒中康复中的应用}

脑卒中是临床上最常见的致残性疾病之一，超过半数幸存者出院时仍存在不同程度的运动功能障碍，其中上肢功能恢复尤其困难。一位卒中后患者重新学会伸手取物转动门把拿起水杯，往往要付出几个月甚至几年的努力。传统康复主要靠治疗师一对一指导，训练强度、训练频次和评估的客观性都受人力资源限制。康复机器人和智能康复系统的引入，为这部分患者带来了新的训练与评估手段。

如图10.12所示，康复机器人在脑卒中康复中的应用通常由患者评估、方案制定、机器人训练、过程数据采集和效果评估反馈等环节构成，多个环节通过数据闭环连接，形成连续动态的康复管理过程。

\begin{figure}[H]
\centering
\begin{tcolorbox}[width=0.88\textwidth, colback=gray!10, colframe=gray!50, boxrule=0.5pt, arc=2pt]
\centering
\vspace{30pt}
{\large\color{gray!70} [图片占位]}
\vspace{10pt}
{\small\color{gray!50} 图10.12 康复机器人在脑卒中康复中的应用流程示意图}
\vspace{25pt}
\end{tcolorbox}
\caption*{\footnotesize\itshape\color{gray!50} 图10.12 康复机器人在脑卒中康复中的应用流程示意图}
\end{figure}

\subsubsection*{\textbf{1. 系统总体架构与训练流程}}

面向脑卒中康复的智能康复系统通常由康复机器人本体、过程采集与控制系统、康复管理平台和远程支持模块构成。康复机器人本体提供训练所需的机械结构、感知与执行能力，常见形式包括末端牵引式上肢训练机器人和外骨骼式上肢训练机器人。过程采集与控制系统负责训练过程中的数据采集、控制策略执行和反馈呈现。康复管理平台用于训练方案制定、训练记录管理、效果评估和团队协作。远程支持模块则在患者出院后或在社区、家庭场景中，提供训练数据上传、专家远程指导和阶段评估等功能。

在训练流程上，系统按五步顺序运行。第一步是患者评估，结合临床检查和功能量表评估患者上肢运动功能、肌力、关节活动度和日常生活能力，确定康复目标和训练强度的初始水平。第二步是方案制定，治疗师在评估结果基础上选择训练任务、设置训练参数并安排阶段性目标。

第三步是机器人训练，按设定方案进行，治疗师在场指导，系统在训练过程中持续采集数据并提供反馈。第四步是过程评估，系统自动汇总训练数据并生成训练报告，由治疗师与医师审阅。第五步是方案调整，根据评估结果对训练任务、参数和阶段目标进行修订，进入下一轮训练。五步循环即构成本节核心的评估、训练、调整康复闭环。

整个训练流程通常以两到四周为一个阶段，不同阶段的训练重点有所不同。早期阶段以恢复关节活动度、唤醒受损通路和建立基本运动模式为主；中期阶段以增强主动控制、提高运动协调性和增加任务难度为主；后期阶段以贴近日常生活功能、提高动作精度和长期保持康复效果为主。智能康复系统可以根据每一阶段的评估结果，对下一阶段的方案进行系统调整。

\subsubsection*{\textbf{2. 患者基线评估与方案制定}}

脑卒中康复的方案制定建立在系统的基线评估之上。常用的功能评估工具包括 Fugl-Meyer 上肢运动功能评估量表、Brunnstrom 分期、改良 Ashworth 痉挛评定量表、徒手肌力测试、关节活动度测量、Wolf 运动功能测试以及 Barthel 日常生活活动能力指数等。

每个量表关注不同的功能维度：Fugl-Meyer 综合评估上肢运动控制能力；Brunnstrom 分期描述运动恢复阶段；改良 Ashworth 评估肌张力变化；Wolf 测试评价上肢功能性动作完成情况；Barthel 指数则评估患者日常生活独立性。它们配合使用，相当于给患者上肢功能做一次立体扫描，是确定训练起点和目标的依据。

基线评估之外，智能康复系统还可结合机器人本体采集的客观数据对患者运动能力进行进一步刻画。例如，通过让患者在机器人辅助下完成若干标准化动作，记录其运动轨迹、关节角度变化、施力大小和稳定性等指标，可以获得对运动控制能力较为细致的描述。这些数据可作为传统量表评估的有益补充，使方案制定建立在主观评估与客观测量相结合的基础之上。

在方案制定过程中，治疗师根据评估结果选择训练任务和参数。对于运动恢复较早期的患者，常以被动训练和辅助训练为主，注重维持关节活动度和激活基本运动模式；对于运动恢复中期的患者，逐步引入主动训练和阻抗训练，强化主动用力和协调控制；对于功能恢复较好的患者，可结合双侧训练、镜像训练和任务导向性训练，进一步贴近日常生活动作。系统会将方案内容、参数设置和训练计划上传至康复管理平台，便于团队成员共同访问和后续调整。

\subsubsection*{\textbf{3. 智能康复训练过程}}

在每次训练过程中，智能康复系统通过传感器实时采集患者的运动轨迹、力的施加情况和肌电信号等信息。系统将运动轨迹与目标轨迹进行实时比较，并通过屏幕上的图形界面向患者呈现轨迹偏差和动作完成度，使患者能够直观了解自身动作质量。施力大小和方向通过力反馈装置传递给患者，使其在主动用力时能够获得真实力感，并在阻抗训练中体验到可调节的负荷。

在控制策略层面，系统通常采用阻抗控制或导纳控制实现“按需辅助”的训练模式。当患者主动用力较强时，机器人减少辅助，让患者主导动作完成；当患者动作出现明显偏差或停滞时，机器人逐步增加辅助力或纠正轨迹，帮助其完成动作。这种交互方式既能够保证训练过程的连续性，又能够最大程度激发患者的主动参与，是智能康复系统区别于传统被动训练设备的重要特征。

为了维持患者训练兴趣，系统通常将训练任务设计为游戏化情境。例如，将上肢伸够动作设计为击落屏幕上的目标、将抓握动作设计为搬运虚拟物体、将精细控制动作设计为绘制特定图形等。游戏化情境配合得分、关卡和成就反馈，可以提高患者的参与度和坚持性。在每次训练结束后，系统会自动生成训练记录，包括完成任务数、动作完成度、用力情况、动作平滑度、轨迹偏差和疲劳变化等指标，并将其同步到康复管理平台，作为后续评估的基础。

\subsubsection*{\textbf{4. 训练效果评估与多维度分析}}

训练效果评估是康复管理中不可缺少的环节。智能康复系统在传统量表评估的基础上，能够提供多维度的客观评估指标。常见评估维度包括运动控制水平、动作完成质量、运动平滑度、轨迹精度、力的稳定性、肌肉激活模式以及训练参与度等。运动控制水平通过任务完成率、目标达到时间和成功完成次数等指标反映；动作完成质量通过轨迹与目标的偏差程度、关节角度变化范围和动作连贯性等指标体现；运动平滑度则可通过加速度变化分析获得，反映动作是否流畅、是否存在不必要的颤动。

智能评估的优势在于能够将原本依赖治疗师主观打分的内容转化为可量化、可比较的客观数据，并在长期训练过程中提供连续记录。例如，对于同一动作，智能康复系统可以在每周训练中对完成时间、轨迹偏差和力的稳定性进行重复测量，并将变化趋势以曲线形式呈现，使治疗师能够更直观地观察患者功能改善情况。多维度评估结果通常以雷达图、趋势曲线和量化报告等形式呈现，便于患者、家属和医务人员共同理解康复进展。

在方法层面，训练效果评估可结合统计分析、机器学习和动作识别技术。统计分析用于刻画指标变化的趋势和差异显著性；机器学习方法可用于建立从客观指标到功能量表评分的映射，使系统在每次训练后自动估计出近似的功能评估分数；动作识别技术用于识别复杂动作中的关键动作要素，并对每一要素的完成情况进行评估。多种方法结合，使训练效果评估能够兼顾整体水平、动作细节和功能意义。

\subsubsection*{\textbf{5. 案例：脑卒中患者上肢康复机器人训练全过程}}

为了更直观地说明系统的应用过程，下面以一例脑卒中后上肢功能障碍患者的康复机器人训练过程为例进行说明。患者男性，58岁，因右侧基底节区脑梗死入院治疗，发病时间约三周，当前生命体征平稳，神志清楚，主要表现为左侧上肢肌力下降、肌张力轻度增高和精细动作能力受限。入院康复评估结果显示，Fugl-Meyer 上肢运动功能评估量表得分为18分（满分66分），Brunnstrom 上肢分期为III期，徒手肌力测试中肩外展和肘伸展约为3级，腕背伸和指伸展约为2级，Barthel 指数为55分。

在患者评估和方案制定基础上，治疗师为其安排了为期4周、每周5次、每次40分钟的上肢康复机器人训练计划。早期一周以辅助式被动训练和肩肘协调训练为主，重点维持关节活动度和激活肩肘运动模式；第二周引入按需辅助式主动训练，通过伸够、抓取虚拟物体等任务提高上肢主动控制能力；第三周开始结合双侧协调训练和轻度阻抗训练，增加任务复杂度；第四周以贴近日常生活的功能性任务为主，包括伸够取物、放置物品和模拟握杯等任务，进一步提升运动协调性和功能可用性。

在4周训练过程中，系统共记录训练数据约 320 万条，包括运动轨迹、关节角度、施力情况和肌电信号等。每次训练结束后，系统自动生成训练日报，每周由治疗师审阅训练周报，每两周进行一次综合评估并对方案进行调整。训练过程中，系统通过动作质量识别功能多次检测到患者使用躯干前倾代偿肩部上抬的现象，并通过提示信息和反馈训练帮助患者逐步纠正；通过疲劳状态评估，系统在多次训练中将中后段难度有所降低，避免出现过度疲劳。

训练结束时的复评结果显示了明显改善：Fugl-Meyer 上肢运动功能评估量表得分由 18 分提升至 32 分，Brunnstrom 上肢分期由 III 期进展至 IV 期，徒手肌力测试中肩外展和肘伸展由 3 级升至 4 级，腕背伸和指伸展由 2 级升至 3 级，Barthel 指数由 55 分提升至 75 分。

系统提供的多维度客观评估结果同样可观：患者在伸够动作中的轨迹偏差由初始平均约 8 厘米下降至约 3 厘米；动作平滑度评分由约 45 分上升至约 70 分；目标完成时间整体缩短约 30\%；伸够、抓握、释放这一组合任务的成功率由约 45\% 提升至约 80\%。这些量化指标，是传统人工评估难以稳定捕捉到的训练进步。

患者出院后继续在社区康复中心和家庭中借助便携设备进行训练，相关训练数据通过远程支持模块持续上传至原住院医院的康复管理平台，由康复团队定期审阅并提供远程指导，脑卒中康复因此从住院期延伸到长期管理期。

\subsubsection*{\textbf{6. 小结}}

康复机器人 + 智能康复系统在脑卒中场景下的价值体现在三处：第一，把训练强度做上去，机器人能稳定提供高重复次数、高一致性的训练；第二，把动作质量看细，多维度客观指标补充了量表的主观评分；第三，把康复时间拉长，从住院期延伸到社区与家庭。58 岁脑梗患者 4 周训练 Fugl-Meyer 18 → 32 分、Barthel 55 → 75 分、动作成功率 45\% → 80\% 的案例，正是医工融合在康复医学中可量化、可推广的典型实例。

\section{中药方剂解读与组方分析}

一位中医师在诊室里，沉吟片刻后写下：桂枝 10g，芍药 10g，生姜 9g，大枣 4 枚，炙甘草 6g。，这就是经典名方桂枝汤。这五味药看似简单，背后却是数千年中医积累出的配伍智慧：哪味药为君、哪味为臣、为什么这样配、加减一味又会变成什么方。面对几万首古今方剂，AI 能不能读懂这套配伍逻辑？这正是中药方剂智能分析要回答的问题。

中医药学是中华民族在长期医疗实践中形成的独立医学体系，方剂学是其重要组成部分。方剂的组方思想包含君臣佐使、性味归经、升降浮沉、寒热补泻等内容，结构丰富、层次复杂。随着古籍数字化、临床病案信息化和药典数据库的建设，方剂相关数据正逐步具备开展人工智能分析的基础，借助自然语言处理、知识图谱和大模型等方法，可以在方剂文本整理、组方分析、辅助阐释和教学应用等方面提供新的支持。

如图10.13所示，中药方剂解读与组方分析通常由方剂文本整理、知识表示与图谱构建、药味关系识别与组方解析以及方剂智能阐释等环节构成，前后衔接形成完整的应用链条。

\begin{figure}[H]
\centering
\begin{tcolorbox}[width=0.88\textwidth, colback=gray!10, colframe=gray!50, boxrule=0.5pt, arc=2pt]
\centering
\vspace{30pt}
{\large\color{gray!70} [图片占位]}
\vspace{10pt}
{\small\color{gray!50} 图10.13 中药方剂解读与组方分析总体框架示意图}
\vspace{25pt}
\end{tcolorbox}
\caption*{\footnotesize\itshape\color{gray!50} 图10.13 中药方剂解读与组方分析总体框架示意图}
\end{figure}

\subsection{方剂文本整理与知识表示}

中药方剂数据散落在古籍、教材、医院病案、药典和方剂数据库等许多地方，同一首桂枝汤，《伤寒论》原文是一种写法，《医方集解》注解是另一种写法，现代《方剂学》教材又是一种写法，医院病案里的桂枝汤加味则是又一种写法。要让人工智能能够理解、检索、比较这些方剂数据，必须先把它们翻译成统一、可计算的形式，这就是本节关注的方剂文本整理与知识表示。

如图10.14所示，方剂文本整理与知识表示通常包括方剂数据采集、文本整理与结构化、知识表示和图谱构建等环节，各环节相互衔接，共同为后续组方分析与智能阐释提供基础。

\begin{figure}[H]
\centering
\begin{tcolorbox}[width=0.88\textwidth, colback=gray!10, colframe=gray!50, boxrule=0.5pt, arc=2pt]
\centering
\vspace{30pt}
{\large\color{gray!70} [图片占位]}
\vspace{10pt}
{\small\color{gray!50} 图10.14 方剂文本整理与知识表示流程示意图}
\vspace{25pt}
\end{tcolorbox}
\caption*{\footnotesize\itshape\color{gray!50} 图10.14 方剂文本整理与知识表示流程示意图}
\end{figure}

\subsubsection*{\textbf{1. 方剂数据的基本特点}}

方剂是中医临床用药的基本组织形式，由两味或两味以上中药按照一定配伍原则组成。一首方剂的档案通常包含方名、组成药味、药物剂量、配伍关系、功效、主治、用法、加减变化和出处等多类信息。

以桂枝汤为例：方名桂枝汤，组成桂枝、芍药、生姜、大枣、炙甘草五味，配伍上以桂枝为君药、配合芍药调和营卫，主治外感风寒表虚证；加减变化多种多样，可衍生出桂枝加葛根汤、桂枝加附子汤等多首方剂。一首方剂往往把上述要素都囊括其中，是较为典型的复杂结构化文本。

方剂数据有四个突出特点。第一，方剂内容由多味药材构成，药味之间的组合并非简单堆砌，而是在君臣佐使等组方思想指导下形成的协同关系。第二，方剂的功效与主治通常以中医理论术语表达，解表散寒、调和营卫、补气养血、清热解毒、益气健脾等，这些术语具有特定的中医含义，需要结合中医基础理论来理解。

第三，方剂内容存在大量表达差异：同一味药在不同文献中可能有多种别名或同义词，如黄芪和黄耆、甘草和炙甘草等；同一首方剂在不同版本中也可能在剂量比例、加减变化和适应证表述上不一致。第四，方剂数据具有较强的历史性和理论性，源自不同时代的方剂在用药习惯、剂量单位和应用情境上各有差异，汉代的一两和现代的克就不能直接画等号。

上述特点决定了方剂数据不能仅按西医文本或一般生物医学数据的方式进行处理，而应在尊重中医药理论体系和表达习惯的前提下，建立适合中医药特点的整理方法和表示方式。这是开展中医药人工智能研究的基本前提，也是保证后续组方分析结果具有医学合理性的重要基础。

\subsubsection*{\textbf{2. 方剂数据来源与采集}}

方剂数据的来源较为广泛，主要包括古籍文献、现代教材、医院病案、药典与标准以及专门方剂数据库等几类。古籍文献是方剂学的重要源头，《伤寒论》《金匮要略》《温病条辨》《医方集解》以及历代本草著作中记载了大量经典方剂，是中医药知识的重要载体。现代教材以《方剂学》《中医内科学》《中医基础理论》等为代表，对古典方剂进行了系统整理，并补充了现代研究成果，是教学和应用中常用的参考资料。

医院病案是方剂在临床实际应用中的重要记录形式，能够反映医师面对具体患者时的方剂选用与加减调整方式，是研究方剂临床应用规律的重要数据来源。药典与标准方面，《中华人民共和国药典》以及相关行业标准对部分常用方剂的组成、剂量、功能主治和质量标准进行了规范化规定，是方剂数据规范化整理的重要依据。

专门的方剂数据库则把上述来源整合后形成结构化资料，常见的包括中国方剂数据库、中医药信息平台、SymMap、TCMID 等，提供方剂、药材、靶点和疾病等多类信息及其相互关系，这些数据库已经成为中医药与人工智能交叉研究的重要数据基础。

在数据采集过程中，需要注意来源差异和版本控制问题。同一首方剂在不同来源中的记录可能存在文字表述、剂量比例和适应证描述等方面的差异；同一来源的不同版本之间也可能存在内容更新与体例差异。为保证后续分析的可靠性，方剂数据采集通常需要建立明确的来源标识、版本记录和适用范围说明，使整理后的方剂数据既能够呈现整体规律，又能够保留来源细节，便于追溯和复核。

\subsubsection*{\textbf{3. 方剂文本的结构化处理}}

方剂文本在原始形式上多为自由文本，包含方名、组成、剂量、用法、功效、主治、加减变化和按语等内容，结构相对松散。要使方剂内容能够被人工智能方法有效处理，需要将其转化为结构化字段，使每首方剂的关键要素能够以统一格式表达。常见结构化字段通常包括方名、出处、组成药味、各药剂量、用法用量、功效、主治证候、适用人群、加减变化和现代研究提示等。

结构化处理通常包括分词与术语识别、字段抽取、术语规范化和关系整理等步骤。分词与术语识别用于在方剂文本中识别出药名、剂量、功效术语和证候术语等基本单元；字段抽取用于将识别结果按照预先设计的字段结构组织起来；术语规范化用于将不同表达形式映射到统一术语，例如将“黄耆”和“黄芪”合并为统一的“黄芪”概念；关系整理用于建立药味与剂量、方剂与功效、方剂与主治之间的对应关系，形成可计算的结构化记录。

在方法层面，方剂文本结构化可以结合规则、机器学习和大模型三种路线。规则方法依靠专家整理的中药术语词典和模板规则，能较准确地处理具有规范表达形式的方剂文本，结果可解释性强；机器学习方法如条件随机场和命名实体识别模型，能在标注语料基础上识别中药名、剂量和功效术语，处理多样化表达。

大模型方法在古文方剂理解、长文本结构化和加减变化分析等任务中表现良好，能够把枳壳与枳实异同这类文献中的细节也讲清楚。但要注意：大模型在医学场景下使用时仍需结合中医药专业词典和规则约束，避免它自由发挥出与古籍记载不符的描述，这是中医药智能分析中尤其重要的边界。

\subsubsection*{\textbf{4. 方剂知识的表示方式}}

方剂知识表示是将整理后的方剂内容以适合计算机处理的形式组织起来的过程。常见表示方式包括结构化表格、三元组、知识图谱和向量表示等。结构化表格以行表示具体方剂、列表示属性字段的方式组织数据，便于数据库存储和基础查询，是较为基础的一种表示方式。三元组表示则将方剂相关知识组织成“主体、关系、客体”的形式，例如“桂枝汤、包含、桂枝”“桂枝、性味、辛温”“四君子汤、主治、脾胃气虚”等，使知识具有更强的关系表达能力。

知识图谱以节点与关系的方式组织方剂相关知识，节点包括方剂、药味、性味、归经、功效、证候和脏腑等概念，关系包括包含、君臣佐使、性味属于、归经于、功效为、主治、配伍和禁忌等。与表格和三元组相比，知识图谱在表达复杂关系和支持图查询方面具有较强优势，能够沿着多跳路径进行检索和推理。例如，根据某一证候节点查找常用方剂，再根据方剂查找其组成药味及性味归经，可以较为便利地获得证候、方剂和药味之间的整体联系。

向量表示是另一种重要的方剂知识表示方式。通过将方剂或药味表示为高维向量，可以使原本以离散符号形式存在的中医药概念具备数值计算属性。常用方法包括基于共现统计的词向量、基于神经网络的嵌入模型以及在知识图谱基础上学习的图嵌入方法。向量表示能够支持相似方剂检索、相似药材推荐和方剂聚类分析等任务，使方剂数据能够被以更细致的方式比较和归纳。在实际应用中，结构化表格、三元组、知识图谱和向量表示常根据任务需要组合使用，由不同方式承担不同分析功能。

\subsubsection*{\textbf{5. 知识图谱构建的基本方法}}

中药方剂知识图谱的构建是方剂智能分析的重要基础。其基本任务是将分散在古籍、教材、药典和数据库中的方剂相关信息整合到统一的图结构中，形成具有较好覆盖度和可扩展性的中医药知识网络。常见构建步骤包括本体设计、实体抽取、关系抽取、知识融合、图谱存储和持续维护。

本体设计用于明确知识图谱中节点和关系的类型体系。中药方剂知识图谱中的节点类型通常包括方剂、药味、性味、归经、功效、证候、症状、脏腑、疾病、剂量单位和用法等；关系类型通常包括“包含”“君臣佐使”“性味为”“归经于”“功效为”“主治”“配伍”“相反”“相畏”和“加减变化为”等。良好的本体设计能够使图谱具有较好的清晰度和扩展性，便于后续不同来源数据的统一接入。

实体抽取与关系抽取从文本中识别出具体的方剂、药材、功效和证候等实体及其相互关系。在中医药场景中，由于古文表达和术语异名较多，实体抽取通常需要结合中药术语词典、方剂名称表和证候分类标准等专业资源；关系抽取则需要结合中医组方理论，对君臣佐使、配伍关系和禁忌关系等进行识别。

知识融合阶段对不同来源的实体和关系进行对齐与归并，例如把《伤寒论》《金匮要略》和现代教材中关于同一首方剂的不同记载整合为统一节点，并在节点上保留各来源标注。图谱存储通常采用图数据库实现，便于支持大规模图查询和路径推理；持续维护包括内容更新、版本管理和质量审核等工作，让图谱内容随着中医药研究和临床实际不断改进。

\subsubsection*{\textbf{6. 小结}}

方剂文本整理与知识表示要做的，是把散落在古籍、教材、病案和药典里的方剂数据翻译成统一可计算的形式。方剂数据具有药味多、关系复杂、理论性强和表达多样的特点；通过分词、术语识别、字段抽取和术语规范化等步骤把自由文本变成结构化记录；再用表格、三元组、知识图谱和向量四种表示方式组合呈现。这一切都建立在尊重中医药理论与表达习惯的前提下，技术再先进，也不能让黄耆和黄芪在图谱中变成两个不相干的节点。

\subsection{药味关系识别与组方解析}

同样的几味药，为什么在不同方剂里配伍出来的效果差异这么大？为什么四君子四味药能成为补气类方剂的核心组合，反复出现在四君子汤、六君子汤、参苓白术散等多首方剂中？AI 能不能从大量方剂数据里看出这些药对和核心组合？这正是本节要回答的问题。

中药方剂的疗效，不仅取决于单味药材的功能，更取决于多味药物之间的相互配合。借助关联规则挖掘、网络分析和大模型语义解析等方法，可以在大量方剂数据中识别药味之间的关系特征，并对方剂结构进行较为系统的解析，把看似相同实则不同的方剂背后的配伍智慧，用可计算的方式呈现出来。

如图10.15所示，药味关系识别与组方解析通常包括药对挖掘、核心药组识别、配伍关系分析、组方模式解析和相似方比较等环节，多个环节相互配合，共同呈现方剂数据中蕴含的配伍规律。

\begin{figure}[H]
\centering
\begin{tcolorbox}[width=0.88\textwidth, colback=gray!10, colframe=gray!50, boxrule=0.5pt, arc=2pt]
\centering
\vspace{30pt}
{\large\color{gray!70} [图片占位]}
\vspace{10pt}
{\small\color{gray!50} 图10.15 药味关系识别与组方解析流程示意图}
\vspace{25pt}
\end{tcolorbox}
\caption*{\footnotesize\itshape\color{gray!50} 图10.15 药味关系识别与组方解析流程示意图}
\end{figure}

\subsubsection*{\textbf{1. 中药药味关系的基本层次}}

中药方剂中的药味关系并非单一层面的简单组合，而是包含多个层次的结构关系。从临床应用角度看，常见层次主要包括药对关系、核心药组关系、组方结构关系和配伍禁忌关系等。药对关系是指在临床使用中频繁出现的两味药组合，如黄芪与当归常用于补气养血、柴胡与黄芩常用于和解少阳、人参与白术常用于补气健脾等。这类两味药的搭配在长期临床实践中形成相对固定的配伍方式，对方剂功效具有较为明确的贡献。

核心药组关系是指由三味或更多药材组成、在多个方剂中反复出现并发挥相对一致功能作用的一组药材。例如，人参、白术、茯苓和炙甘草共同构成的“四君子”药组，是补气健脾类方剂的核心组合，常出现在四君子汤、六君子汤、香砂六君子汤、参苓白术散等多首方剂中。组方结构关系是指方剂内部在君臣佐使、寒热并用、补泻兼施等组方思想指导下形成的整体结构关系，反映了方剂的整体设计思想。

配伍禁忌关系是中药方剂中需要特别关注的内容。中医药传统经验中归纳了“十八反”“十九畏”等配伍禁忌，描述了部分药材在同时使用时可能出现的不良影响或药效减弱情况。在方剂智能分析中，应当将配伍禁忌作为重要的约束条件加以考虑，避免将禁忌组合误识别为合理配伍。理解药味关系的层次性，对正确开展药味关系识别和组方解析具有重要意义。

\subsubsection*{\textbf{2. 药对与高频共现挖掘}}

药对挖掘是药味关系识别中最基础的任务之一，主要依托关联规则挖掘和频繁项集分析等方法。其基本思路是将每首方剂视为一个“事务”，方剂中包含的药材视为该事务中的“项”，通过统计不同药材在方剂中的共现频率和关联强度，挖掘出在方剂数据集中高频出现的药对。常见挖掘算法包括 Apriori 算法、FP-Growth 算法等，能够在大规模方剂数据中较高效地发现频繁共现的药材组合。

在方剂数据分析中，常用统计指标包括支持度、置信度和提升度等。支持度反映某一药对在所有方剂中出现的比例，用于衡量该药对在整体方剂数据中的普遍程度；置信度反映在出现某一味药的方剂中同时出现另一味药的比例，用于衡量两味药之间的关联强度；提升度则用于在排除随机共现影响的基础上评价两味药之间是否存在显著关联。综合使用三类指标，可以较准确地识别出真正具有配伍意义的药对，避免将偶然共现误认为稳定关系。

药对挖掘的结果通常需要结合中医药专业知识进行解读。某一药对在数据中虽然出现频率较高，但其实际配伍意义仍需结合性味、归经、功效和主治证候加以理解。例如，某些药对在多种方剂中共同出现，可能反映特定证候下的常规配伍方式，也可能与某一时期的用药习惯或文献来源结构有关。在实际研究中，应当将算法挖掘结果与中医专家解读相结合，使数据分析既具有数量上的支持，又能够保持医学合理性。

\subsubsection*{\textbf{3. 网络分析与组方模式识别}}

在药对挖掘基础上，可以进一步借助网络分析方法对方剂数据进行整体研究。其基本思路是将每味药材视为网络中的节点，将药材之间的共现关系视为网络中的连边，连边的权重通常与药对共现频率或关联强度相对应。通过网络构建，可以将分散的方剂数据整合为一张综合的“药材关系网络”，便于从全局角度观察药材之间的相互关系。

在网络分析中，常用方法包括节点中心性分析、社区发现、子图模式识别和网络可视化等。节点中心性分析用于识别在网络中处于中心位置的关键药材，它们与多种其他药材广泛配伍，常常是某一类治法中的主力药味。社区发现用于识别网络中联系紧密、相对独立的药材社区，每一社区往往对应某一类证候或某一种治法的核心药组。

子图模式识别用于发现具有特定结构特征的药材组合，比如稳定的三味药或四味药核心组合；网络可视化则把分析结果以图形方式呈现，让研究者一眼就能把握药材关系的整体结构。把方剂数据画成一张网，正是这种分析方法最直观的优势。

网络分析的优势在于能够在较大规模数据中发现潜在的组方模式。例如，对补气类方剂数据进行网络分析时，常可观察到人参、黄芪、白术、茯苓和炙甘草等药材形成紧密的核心区，体现了补气健脾类方剂的常见配伍特征；对清热解毒类方剂进行分析时，则可观察到金银花、连翘、黄芩、栀子和板蓝根等药材的紧密关联。这些网络结构特征与中医临床经验之间往往存在较好对应，可作为辅助理解组方规律的重要参考。

\subsubsection*{\textbf{4. 君臣佐使与组方功能解析}}

君臣佐使是中医方剂学中的重要组方思想，描述了方剂内部药味之间的功能层次关系。君药通常是方剂中针对主病或主证起主要治疗作用的药味，体现方剂的核心治法；臣药辅助君药增强主治作用，或针对兼证发挥治疗作用；佐药用于协助君臣两类药材，增强药效、缓和峻烈或针对次要兼证；使药则用于引导诸药直达病所，或调和诸药使全方协同发挥作用。理解和识别方剂中的君臣佐使结构，是开展方剂智能解析的关键内容之一。

在方剂智能解析中，可以借助多种方法对君臣佐使关系进行识别和验证。基于规则的方法依据剂量大小、药材性味、功效与主治证候的对应关系等条件，对方剂中各药材在方剂中的角色进行初步判断；基于知识图谱的方法将药材性味、归经、功效和方剂主治等信息组合起来，依据图谱中的关系路径判断各药材与方剂主治证候之间的关联紧密程度；基于大模型的方法则依托对古籍和教材文本的语义理解，对方剂组方思想进行整体分析。

在组方功能解析方面，可结合方剂主治证候和组成药材的功效特点，识别方剂的主攻方向。例如，补气健脾类方剂以人参、白术、茯苓、炙甘草等补气药材为主，配合少量行气或祛湿药材，主治气虚和脾胃虚弱相关证候；清热解毒类方剂以金银花、连翘、黄芩、栀子等清热药材为主，主治热毒壅盛相关证候；活血化瘀类方剂以丹参、川芎、赤芍、桃仁等活血药材为主，主治瘀血相关证候。

通过组方功能解析，方剂的整体设计思想就能以更直观的形式呈现：一首方剂主攻什么、靠哪几味药主攻、辅以什么协同，这正是中医师选方加减时的核心思路，也是方剂学习者最希望掌握的技能。

\subsubsection*{\textbf{5. 相似方比较与方剂家族}}

在中医方剂学习和应用中，常会遇到组成相近、功能相似的方剂。这些方剂之间往往存在历史源流关系，通过加减药味或调整剂量从基础方剂衍生而来，构成所谓的“方剂家族”。例如，由四君子汤为基础，加陈皮形成异功散，加陈皮和半夏形成六君子汤，再加木香和砂仁形成香砂六君子汤等，这些方剂在组成上具有较为明显的层次关系。借助人工智能方法对相似方进行比较，有助于学习者把握方剂之间的演变规律和应用差别。

相似方比较通常包括组成对比、剂量比例对比、功效与主治对比和适用情境对比等内容。组成对比用于识别两首方剂之间共有的药味和各自独有的药味，揭示加减变化的具体方向；剂量比例对比用于分析在共有药味中各方剂的剂量差异，反映方剂针对的证候差别；功效与主治对比用于比较两首方剂在治疗重点和适用证候上的不同；适用情境对比则结合临床应用背景，指出在何种情况下应优先选用哪一方剂。

在方法层面，相似方比较可结合知识图谱、向量表示和大模型解析等多种方法。知识图谱提供方剂之间组成、功效和主治的结构化对照；向量表示可在嵌入空间中度量方剂相似度，并支持相似方推荐和方剂聚类；大模型则在比较结果基础上生成自然语言形式的解读，帮助学习者更直观地理解方剂之间的差异。

需要特别指出：方剂之间的相似度计算结果不能直接等同于临床等价关系。两首方剂在数据上看似接近，并不意味着可以在临床上随意替换，患者的体质、合并症和用药史都可能影响选方。AI 给出的相似度是一种参考，最终的临床决策仍由中医师作出。

\subsubsection*{\textbf{6. 小结}}

药味关系识别与组方解析回答的是中医方剂学的核心问题：方剂里的药味怎么配合、为什么这样配。从最基础的药对挖掘（如黄芪、当归、柴胡、黄芩），到核心药组识别（如四君子组合贯穿多首方剂），再到借助网络分析揭示组方模式、用君臣佐使解析方剂主攻方向、通过相似方比较把握方剂家族演变，这些方法把传统中医配伍智慧用算法表达了出来。但所有结果都要结合中医药专业知识解读，把数据支持和理论合理性结合在一起，才能为后续方剂智能阐释和教学应用提供真正可信的基础。

\subsection{案例：基于知识图谱的四君子汤组方解读}

前两节我们准备好了方剂的知识底座，把古籍、教材、病案里的方剂数据变成结构化记录和知识图谱，把药味关系、君臣佐使和方剂家族也分析了出来。本节把这些零件组装成一个面向用户的应用：当一位学生、一位中医师或一位对中医感兴趣的普通读者在屏幕上输入四君子汤三个字，系统能给出一份怎样的解读？我们就以这首最经典的补气方为例，把AI 读方剂的过程完整呈现一遍。

如图10.16所示，基于知识图谱的方剂组方解读系统通常由方剂输入、知识图谱检索、大模型语义生成、解读结果输出和反馈学习等环节构成，多个环节通过统一的数据接口连接，形成端到端的方剂解读流程。

\begin{figure}[H]
\centering
\begin{tcolorbox}[width=0.88\textwidth, colback=gray!10, colframe=gray!50, boxrule=0.5pt, arc=2pt]
\centering
\vspace{30pt}
{\large\color{gray!70} [图片占位]}
\vspace{10pt}
{\small\color{gray!50} 图10.16 基于知识图谱的方剂组方解读系统总体架构示意图}
\vspace{25pt}
\end{tcolorbox}
\caption*{\footnotesize\itshape\color{gray!50} 图10.16 基于知识图谱的方剂组方解读系统总体架构示意图}
\end{figure}

\subsubsection*{\textbf{1. 系统总体架构与解读流程}}

面向方剂智能阐释的系统通常由数据层、知识层、模型层和应用层四部分组成。数据层负责存储方剂原始资料，包括古籍文献、现代教材、药典记载和病案资料等；知识层以中药方剂知识图谱为核心，将方剂、药味、性味、归经、功效、证候和加减变化等内容组织为结构化网络；模型层以大模型为核心，结合检索增强生成机制将知识图谱内容用于约束生成结果；应用层面向教学、培训、科普和临床查阅等不同场景，以可视化界面和结构化报告形式呈现解读结果。

在解读流程上，系统按五步顺序运行。第一步是方剂输入，用户通过名称查询或药味输入指定需要解读的方剂。第二步是知识图谱检索，系统在图谱中查找该方剂的组成、剂量、性味归经、功效、主治证候、衍生方剂和使用注意等信息。第三步是大模型语义生成，模型在检索结果基础上生成自然语言形式的解读内容。

第四步是解读结果输出，系统把生成内容与图谱信息整合为结构化报告，呈现给用户。第五步是反馈学习，使用者对解读结果的评价和修正反馈到系统，用于规则优化和图谱完善。这五步走完，一份方剂的立体说明书就出现在用户屏幕上。

整个解读流程通常在十几秒至数十秒内完成。系统会将解读内容按照组成与剂量、君臣佐使、功效主治、适用人群与使用注意以及衍生方剂等模块组织，使用户可以根据需要展开查看具体内容。系统还会在必要位置标注信息来源，例如某一表述来自古籍原文、某一现代认识来自药典或学术研究，使解读结果具有较好的可追溯性。

\subsubsection*{\textbf{2. 知识图谱与大模型的协同机制}}

在方剂阐释系统中，知识图谱与大模型通常以协同方式工作。知识图谱负责提供方剂相关的事实性内容，包括方剂组成、药味性味、组方结构和主治证候等信息；大模型负责将这些内容组织成具有可读性的自然语言表达，并对组方思想进行综合性解读。两者结合后，能够在保证内容准确性的前提下提升解读结果的清晰度和友好程度。

在具体技术路线上，系统采用检索增强生成（RAG）机制。当用户提交方剂查询请求时，系统首先在知识图谱中检索与该方剂相关的节点和关系，获得组成药味、剂量、性味、归经、功效和主治证候等结构化信息；随后把这些信息以提示词的形式输入大模型，让它在限定知识范围内生成解读内容。这就像让大模型在开卷考试中作答，既发挥了它在语言表达上的优势，又通过知识图谱对生成内容做了事实层面的约束，降低了自由发挥出与中医药事实不符内容的风险。

在协同机制中，知识图谱还承担事实校验功能。大模型生成内容后，系统可以将关键信息与图谱内容进行比对，识别可能存在的偏差并加以提示。例如，若生成结果在药味组成、性味归经、功效或主治证候等方面与图谱内容存在差异，系统会自动标注或提示，使最终呈现的内容能够与图谱事实保持一致。这种“图谱约束、模型生成、图谱校验”的循环机制，是方剂智能阐释系统在中医药场景中的重要保障。

\subsubsection*{\textbf{3. 方剂多维度解读维度}}

方剂智能阐释通常涉及多个解读维度。第一是方剂组成与剂量维度，介绍方剂中各味药材的名称、剂量和用法，必要时附原方出处和剂量换算说明，使用户能够快速了解方剂的基本组成。第二是君臣佐使维度，分析各味药材在方剂中的地位与作用，结合药材性味、归经和功效特点说明组方思想。第三是功效主治维度，对方剂的整体功效和主治证候进行解读，包括传统中医表述和较为通俗的语言说明。

第四是适用人群与使用注意维度，介绍方剂适合的人群特征、不宜使用的人群、常见使用注意以及与其他药物或食物之间的注意事项。第五是衍生方剂维度，介绍由该方剂加减变化形成的相关方剂及其应用差别，使用户能够把握方剂之间的演变关系。第六是现代药理研究提示维度，介绍与该方剂相关的现代药理研究、临床应用研究和质量标准内容，便于用户了解方剂在现代研究中的进展。

不同解读维度面向不同应用需求。教学场景中通常需要完整呈现各类维度内容，便于学生系统理解；临床查阅场景中常重点关注组成与剂量、功效主治和使用注意等维度，便于快速参考；科普场景则更强调通俗解读和使用注意，避免过多专业术语。系统通常根据应用场景和用户角色，对不同维度的内容深度和呈现方式进行调整，使解读结果与具体使用情境相匹配。

\subsubsection*{\textbf{4. 案例：四君子汤的组方解读全过程}}

为了更直观地说明系统的解读过程，下面以四君子汤为例进行介绍。四君子汤出自《太平惠民和剂局方》，由人参、白术、茯苓和炙甘草四味药组成，是补气健脾类方剂的代表性方剂之一。系统接收到“四君子汤”的查询请求后，约 12 秒生成完整的解读报告，按组成与剂量、君臣佐使、功效主治、适用人群与使用注意以及衍生方剂等模块呈现内容。

在组成与剂量模块中，系统列出四君子汤的传统组成：人参约 9 克，白术约 9 克，茯苓约 9 克，炙甘草约 6 克，水煎服。系统在介绍组成时同时给出方剂出处和原方剂量记载，并指出在临床实际应用中可根据患者具体情况进行剂量调整。

在君臣佐使模块中，系统通过知识图谱检索结合大模型语义生成给出分析结果。人参为君药，性味甘温，主入脾肺经，具有大补元气、补脾益肺的作用，是方剂中针对气虚主证的核心药味；白术为臣药，性味甘苦温，主入脾胃经，具有健脾益气、燥湿利水作用，辅助君药增强补气健脾之力。

茯苓为佐药，性味甘淡平，主入心脾肺经，具有健脾渗湿作用，配合白术运化水湿，使补气而不滞腻；炙甘草为使药，性味甘平，主入心肺脾胃经，具有补脾益气、调和诸药的作用，使全方协同发挥作用。一首方剂的君臣佐使关系就这样以可解释、可追溯的方式呈现在用户面前。

在功效主治模块中，系统把四君子汤的传统功效概括为补气健脾、益气养胃，主治脾胃气虚证，常见表现为面色萎黄、语声低微、四肢无力、食欲不振、大便溏薄、舌淡苔白、脉虚弱等。系统在专业表述之外，还以较为通俗的语言指出：该方剂主要用于因脾胃功能较弱引起的气短乏力、胃口不好、容易疲劳等情况，是中医补气健脾类方剂的基础方。

在适用人群与使用注意模块中，系统说明四君子汤适合脾胃气虚体质者，对于阴虚火旺、湿热壅盛或外感发热等情况则不宜单独使用；在合并其他疾病、正在服用其他药物或处于特殊生理状态时，应在中医师指导下使用，系统在这一模块明确把知识参考和个体诊疗区分开来。

在衍生方剂模块中，系统利用知识图谱中方剂之间的加减关系，列出四君子汤的常见衍生方剂及其差别。异功散在四君子汤基础上加陈皮，主治脾胃气虚兼气滞，可表现为腹胀、嗳气等症状；六君子汤再加半夏，主治脾胃气虚兼痰湿，可表现为咳嗽痰多、胸闷等症状。

香砂六君子汤在六君子汤基础上加木香、砂仁，主治脾胃气虚兼气滞痰湿、食欲不振；参苓白术散则在四君子汤基础上加山药、扁豆、莲子、薏苡仁、砂仁、桔梗等，主治脾虚湿盛、泄泻便溏。这条由四君子出发的方剂家族演变线，被系统以结构化表格和文字说明相结合的方式呈现，让学习者一目了然地把握方剂之间的传承与差异。

在解读结果输出后，系统会请用户根据使用情况对内容进行评价。教学使用者可对君臣佐使分析的清晰程度、衍生方剂介绍的准确性以及通俗解读的友好程度进行打分，临床使用者可对组成剂量、使用注意和现代药理研究提示的实用性进行评价，反馈数据汇总后用于规则优化和图谱完善。系统还会保留每次解读结果的版本记录和数据来源，便于在内容更新和版本管理过程中进行追溯。

\subsubsection*{\textbf{5. 应用场景与伦理边界}}

基于知识图谱的方剂组方解读系统在中医药教育、医师培训、临床查阅和大众科普等场景中均具有应用价值。在中医药教学中，系统可以帮助学生系统理解经典方剂的组成、组方思想和加减变化，并通过可视化的图谱关系增强学习兴趣。在中医师规范化培训中，系统可作为临床查阅工具，帮助住院医师和专科医师在面对具体方剂时快速获得规范化的参考信息。在面向大众的中医药科普中，系统可以以通俗易懂的方式介绍常见方剂的功效与使用注意，引导大众形成科学认识。

在应用过程中，方剂智能阐释的伦理边界需要特别强调。第一，方剂解读结果是知识性参考，不应被视为针对个体患者的诊疗建议，使用者不应依据系统输出自行调整用药方案。第二，中医药临床应用强调辨证施治，方剂使用需要由中医师根据患者具体情况开具，相关诊疗活动应由具有资质的中医师承担，智能阐释系统不能替代中医师的临证判断。

第三，中医药文化具有深厚的历史底蕴，对古籍内容的引用应当保持准确，避免出现伪中医内容传播或文化失真。系统设计和使用过程中，应当通过数据来源标注、内容审核机制和使用界面提示等方式明确这些边界，使方剂智能阐释在合理范围内发挥作用，这既是技术伦理的要求，也是对中医药文化的尊重。

\subsubsection*{\textbf{6. 小结}}

基于知识图谱的方剂组方解读系统，把方剂的事实骨架（知识图谱）和自然语言能力（大模型）拼到了一起：图谱保证准确，大模型保证可读，图谱约束、模型生成、图谱校验循环让两者协同。四君子汤案例显示，AI 可以从组成、君臣佐使、功效主治到衍生方剂多个维度给出可读且专业一致的阐释。前提是清楚的边界意识，AI 是知识参考，不是临证医师；尊重中医师的临证主导权和中医药文化底蕴，方剂智能阐释才能真正服务于教学、培训和科普。

\section*{本章小结}

本章把视野从AI 看病历影像扩展到AI 走进医工融合的具体场景，沿四个方向展开：医嘱智能判别与规范评分、智能穿戴监测与健康评估、智能康复系统与功能训练，以及中药方剂解读与组方分析。它们的共同点在于：人工智能不再是单点的算法 demo，而是要与医学知识、工程系统、临床流程深度融合，既深入理解临床业务，又充分发挥工程方法和算法模型的能力。

医嘱智能判别与规范评分把几秒钟内审一份医嘱做成了可能。从医嘱要素拆解、结构化抽取，到金标准医嘱比对、规则提取，再到知识图谱构建，最后由大模型与图谱协同生成医嘱体检报告，AI 在保留医师主导权的前提下，识别检查项目缺失、用药不规范、监测安排不足、配伍禁忌和方案与指南偏差等问题，为医嘱质量管理提供了有效辅助。

智能穿戴监测与健康评估让健康观察从间断照片变成长视频。可穿戴设备把心率、血压、血糖、心电、活动和睡眠数据从医院延伸到院外；健康状态识别通过特征提取、模型识别和个体化基线读懂这些数据；分级预警与人机协同干预把异常发现转化为合理的健康行动。72 岁慢性心衰患者 90 天 3 次预警 / 2 次方案调整 / 1 次避免再入院的案例，展示了这套系统在慢性病管理和老年人健康保障中的真实价值。

智能康复系统与功能训练把传统康复器械的只出力升级为会感知、能配合。以感知、执行、控制、交互四个模块为骨架，结合被动 / 主动 / 辅助 / 阻抗多种交互方式和阻抗 / 导纳 / 混合等控制策略，AI 可以做到动作质量识别、代偿动作识别、疲劳评估和难度自适应调整，并通过日报、周报、阶段评估的数据闭环让训练可追溯。58 岁脑梗患者 4 周训练 Fugl-Meyer 18→32 分、Barthel 55→75 分的案例，是康复机器人 + 智能康复系统价值的具体写照。

中药方剂解读与组方分析把数千年中医的配伍智慧翻译成可计算的形式。方剂文本整理与知识表示把古籍、教材、病案和药典里的数据组织成结构化网络；药味关系识别与组方解析借助关联规则挖掘、网络分析和大模型语义解析揭示药对、核心药组、君臣佐使结构和方剂家族；基于知识图谱的方剂解读系统则把事实骨架和自然语言能力拼到一起，对四君子汤这类经典方剂给出可读且专业一致的阐释。在所有环节中，尊重中医药理论体系和中医师的临证主导权始终是底线。

通过本章学习可以看到：医工融合方向的人工智能应用，从来不是单点算法的展示，而是医学知识、工程方法和算法模型在真实场景下的协同，医嘱审核背后是临床指南、知识图谱与大模型的合力；穿戴监测背后是传感器、云端、医护协同的接力赛；康复机器人背后是机械、控制与神经可塑性理论的结合；方剂解读背后是中医药文化与现代知识工程的对话。AI 是辅助工具，不是替代者，这是医工融合应用真正能够长久发展的前提。

\section*{课后习题10}

\begin{enumerate}
\item 医嘱由哪四类基本要素构成？请用一句话概括每类要素回答的核心问题，并说明其中任一要素缺失会带来什么影响。

\item 金标准医嘱和规则提取在医嘱智能审核中各自承担什么作用？请结合前提、结论形式的规则举一两个例子加以说明。

\item 知识图谱与大模型在医嘱智能审核中如何协同？请用开卷考试或类似的比喻向一位非计算机专业的同学解释 RAG 机制。

\item 与院内常规测量相比，可穿戴设备的连续监测在临床发现能力上有什么独特价值？请举出阵发性房颤、夜间高血压或糖尿病血糖波动中的至少两个具体例子。

\item 健康状态识别中常用的异常识别方法有哪三类？以一位运动员的静息心率为例，说明个体化基线建模为什么是必要的。

\item 智能穿戴监测系统中分级预警机制为什么要分为轻度提示、中度预警和重度告警三级？过于宽松或过于敏感各会带来哪些问题？

\item 智能康复设备由哪些硬件模块构成？常见的四种人机交互方式（被动、主动、辅助、阻抗）分别适合康复的哪个阶段？请举例说明。

\item 代偿动作识别在康复训练中为什么重要？请结合上肢训练中以躯干前倾代替肩部上抬的例子，说明若不及时纠正可能造成的后果。

\item 中药方剂数据具有哪些独特特点？为什么不能直接照搬西医文本或一般生物医学数据的处理方式？请结合黄耆与黄芪等具体例子说明。

\item 关联规则挖掘和网络分析在方剂数据上分别能发现哪类组方规律？请以四君子核心药组或补气类方剂的常见组合为例加以说明。

\item 基于知识图谱的方剂组方解读系统在使用过程中应当注意哪些伦理边界？为什么AI不能替代中医师作出临证决策？请结合中医药文化与诊疗特点谈谈你的理解。

\end{enumerate}
